\chapter{The Road to 6G}
\label{ch:sixg}

\begin{chapterintro}
Every chapter so far has described technology that exists: standards that are written, chips that
ship, links that carry traffic today. This last chapter is about the next one. A sixth generation of
mobile systems is being defined as this book is written, with the ITU's framework for
``IMT-2030'' agreed in 2023, 3GPP's first 6G studies under way in Release~20 and the first
specifications planned for Release~21 around the end of the decade. The chapter does two jobs. First,
it explains, with the same rigour as the rest of the book, the physical-layer ideas that the 6G
debate revolves around: the upper mid-band and sub-terahertz spectrum and the hardware limits there;
delay--Doppler waveforms (OTFS and AFDM) derived from first principles; near-field beamfocusing with
very large arrays, cell-free massive MIMO and reconfigurable intelligent surfaces; integrated sensing
and communication; learned air interfaces; the integration of satellites and high-altitude platforms;
semantic communication, energy efficiency, zero-energy devices, post-quantum security, quantum key
distribution and optical wireless. Second, it tries to separate physics from fashion. Some of these
ideas will be in your phone in 2032; some will not survive contact with a power amplifier, a
regulator or a business case. The chapter ends with a section of a different kind: what, after
twenty-five chapters, is worth carrying with you for the rest of a career.
\end{chapterintro}

%======================================================================
\section{Why another generation?}
\label{sec:ch25:why}
%======================================================================

A generation of mobile technology is not a law of nature. It is a bundle of three things that
happen to move together roughly once a decade: a new radio interface that is not backward compatible
with the old one, new spectrum to deploy it in, and a new business story that justifies the
investment. Chapter~\ref{ch:cellular} told the story from AMPS to 5G. Each transition was driven by
a limit the previous generation could not escape: analog FM could not be encrypted or packed tightly
enough, GSM's circuit-switched TDMA could not carry packet data efficiently, WCDMA's single carrier
could not scale to 20\,MHz and MIMO gracefully, and LTE's always-on reference signals, fixed numerology
and lack of a millimetre-wave design made it a poor fit for the ultra-reliable, massive-IoT and
multi-gigabit services that 5G promised.

Is there such a limit for 5G? The honest answer is ``fewer than before''. 5G NR is a flexible
framework: scalable numerology, beam-based operation from 410\,MHz to 71\,GHz, a lean carrier, LDPC
and polar codes within a fraction of a decibel of capacity, massive MIMO, and a service-based core.
It can, and through the 5G-Advanced releases 18 and 19 does, absorb many of the ideas proposed for 6G.
The case for a new generation therefore rests on four arguments, which recur throughout this chapter:
\begin{enumerate}
\item \textbf{New spectrum needs new design.} Capacity is still, to first order, bandwidth times
spectral efficiency times the number of cells (Chapter~\ref{ch:multipleaccess}). The bands below
7\,GHz are full. The next large allocations are in the \emph{upper mid-band}, 7--24\,GHz, where
arrays of hundreds of elements become practical, and further out in the sub-terahertz range.
\item \textbf{The device count and diversity keep growing.} Machines, sensors and vehicles have very
different needs from smartphones, and some of them have no battery at all.
\item \textbf{The network can do more than move bits.} A base station already measures delay,
Doppler and angle of every user; turning it into a radar, a positioning system or a source of
training data for machine learning is a natural next step.
\item \textbf{Cost and energy dominate operator economics.} A new generation is a chance to remove
legacy overhead, to make power consumption proportional to traffic, and to simplify.
\end{enumerate}

\begin{history}[The decade rhythm]
The first commercial cellular networks opened in Tokyo (NTT, 1979), in the Nordic countries (NMT,
1981) and in Chicago (AMPS, 1983). The first GSM call on a commercial network was made in Finland in
1991. NTT DoCoMo launched the first WCDMA service, FOMA, in October 2001; TeliaSonera opened the first
commercial LTE networks in Stockholm and Oslo in December 2009; and the first 5G smartphone services
began in South Korea and the United States in April 2019. The rhythm of roughly one generation per
decade is not an accident. It matches the time needed to agree a standard, to clear and auction
spectrum, to amortise the previous network, and to move an integrated circuit from a research
prototype to a handset. Every forecast of 6G arriving ``around 2030'' is really a forecast that this
rhythm will hold once more.
\end{history}

\begin{keyidea}[A generation is a set of bets]
Read every 6G proposal as a bet: that a new waveform will pay for its receiver complexity, that a
new band can be cleared and covered economically, that a passive surface is cheaper than a small
cell, that a neural network will generalise outside its training set. The physics in this chapter
lets you estimate the odds. The standards process and the market settle them.
\end{keyidea}

%======================================================================
\section{The process: IMT-2030 and 3GPP}
\label{sec:ch25:process}
%======================================================================

\subsection{ITU-R and the IMT-2030 framework}

Formally, a ``generation'' is defined by the ITU Radiocommunication Sector (ITU-R). Its Working
Party 5D, which produced the IMT-2000 (3G), IMT-Advanced (4G) and IMT-2020 (5G) families, adopted
Recommendation \textbf{ITU-R M.2160}, ``Framework and overall objectives of the future development of
IMT for 2030 and beyond'', in late 2023. The name of the new family is \term{IMT-2030}. The
Recommendation does not specify any technology. It describes six \emph{usage scenarios}, four
\emph{overarching aspects} that apply to all of them, and fifteen \emph{capabilities}, with example
target values ``for research and investigation''.

\begin{figure}[htbp]\centering
\begin{tikzpicture}[
  sc/.style={draw=navy,fill=navy!8,thick,regular polygon,regular polygon sides=6,minimum size=2.5cm,align=center,font=\sffamily\scriptsize,inner sep=0pt},
  nw/.style={draw=accent,fill=accent!8,thick,regular polygon,regular polygon sides=6,minimum size=2.5cm,align=center,font=\sffamily\scriptsize,inner sep=0pt}]
\node[sc] (a) at (90:2.7)  {Immersive\\communication\\{\tiny (from eMBB)}};
\node[sc] (b) at (30:2.7)  {Hyper-reliable\\low-latency\\{\tiny (from URLLC)}};
\node[sc] (c) at (-30:2.7) {Massive\\communication\\{\tiny (from mMTC)}};
\node[nw] (d) at (-90:2.7) {Ubiquitous\\connectivity};
\node[nw] (e) at (-150:2.7){Integrated\\sensing and\\communication};
\node[nw] (f) at (150:2.7) {AI and\\communication};
\node[draw=navy!60,fill=white,circle,minimum size=2.3cm,align=center,font=\sffamily\tiny,text width=1.6cm,inner sep=1pt] at (0,0)
 {\textbf{overarching}\\sustainability\\connecting the\\unconnected\\ubiquitous\\intelligence\\security and\\resilience};
\end{tikzpicture}
\caption{The six usage scenarios of IMT-2030 (Recommendation ITU-R M.2160). The three in navy extend
the eMBB, URLLC and mMTC scenarios of IMT-2020; the three in red are new. The four overarching
aspects in the centre apply to every scenario.}
\label{fig:ch25_scenarios}
\end{figure}

Figure~\ref{fig:ch25_scenarios} shows the scenarios. Three of them extend the familiar 5G triangle:
\emph{immersive communication} (extended reality, holographic-style telepresence, high-rate mobile
broadband), \emph{hyper-reliable and low-latency communication} (industrial control, telemedicine,
coordinated vehicles) and \emph{massive communication} (dense sensor and device populations). Three
are new. \emph{Ubiquitous connectivity} targets the users that terrestrial networks do not reach, and
is the hook for satellites and high-altitude platforms. \emph{AI and communication} covers both the
use of AI inside the network and the network as a platform for distributed learning and inference.
\emph{Integrated sensing and communication} makes sensing (detecting, locating, tracking and imaging
objects that carry no device) a service of the mobile network.

The capabilities are a mixture of the old and the new. Nine are enhancements of IMT-2020 capabilities
(peak and user-experienced data rate, spectral efficiency, area traffic capacity, connection density,
mobility, latency, reliability, and security, privacy and resilience). Six are new: coverage,
positioning, sensing-related capabilities, AI-related capabilities, sustainability and
interoperability. Table~\ref{tab:ch25:kpi} lists the example values given in M.2160 next to the
IMT-2020 minimum requirements of Report ITU-R M.2410, and Figure~\ref{fig:ch25_kpi} plots the
ratios.

\begin{table}[htbp]\centering\small
\caption{Selected IMT-2030 capabilities: example values for research and investigation from
Recommendation ITU-R M.2160, compared with the IMT-2020 minimum technical requirements (Report ITU-R
M.2410). The IMT-2030 values are not requirements; the formal requirements are being developed by
WP~5D for the evaluation phase.}
\label{tab:ch25:kpi}
\begin{tabularx}{\linewidth}{@{}X l l@{}}
\toprule
Capability & IMT-2020 requirement & IMT-2030 example values\\
\midrule
Peak data rate (downlink) & 20\,Gb/s & 50, 100, 200\,Gb/s\\
User-experienced data rate & 100\,Mb/s & 300, 500\,Mb/s\\
Spectrum efficiency & 30\,b/s/Hz peak (DL); averages per environment & 1.5 and 3 times IMT-2020\\
Area traffic capacity & 10\,Mb/s/m$^2$ & 30, 50\,Mb/s/m$^2$\\
Connection density & $10^6$ devices/km$^2$ & $10^6$--$10^8$ devices/km$^2$\\
Mobility & 500\,km/h & 500--1000\,km/h\\
User-plane latency & 1\,ms (URLLC) & 0.1--1\,ms\\
Reliability & $1-10^{-5}$ & $1-10^{-5}$ to $1-10^{-7}$\\
Positioning accuracy & (not specified) & 1--10\,cm\\
Coverage, sensing, AI, sustainability & --- & qualitative; to be defined\\
\bottomrule
\end{tabularx}
\end{table}

\mdcfig[0.88]{ch25_kpi}{IMT-2030 example values (red bars, the range between the lowest and highest
example in Recommendation ITU-R M.2160) relative to the IMT-2020 minimum requirements (dots at 1).
Apart from connection density and reliability, the targets ask for factors of 1.5 to 10, not the
factors of 10 to 100 that separated earlier generations. The headline ambitions of 6G lie mostly in
the new capabilities---sensing, positioning, AI, sustainability---that have no 5G baseline.}

Two things stand out. First, the improvement factors are modest. The jump from 4G to 5G was a factor of
20 in peak rate and a factor of ten in latency; the 6G examples ask for a few times more. The research
literature of 2019--2021 that promised ``1\,Tb/s and 0.1\,ms everywhere'' has been tempered by the
economics of 5G deployment. Second, the new capabilities are qualitative. How should the ``sensing
capability'' of a network be specified and tested? Nobody yet knows, and that is part of the work.

\subsection{The timeline}

The ITU process (Figure~\ref{fig:ch25_timeline}) follows the pattern of IMT-2020: the vision
(M.2160, 2023); a report on technical performance requirements and the evaluation methodology
(scheduled for completion around 2026); a call for candidate technologies with submissions and
independent evaluation in roughly 2027--2029; and a Recommendation specifying the IMT-2030 radio
interfaces around 2030. The World Radiocommunication Conferences frame the spectrum: WRC-23 placed
on the agenda of \textbf{WRC-27} (agenda item~1.7) studies of possible IMT identifications in
4.4--4.8\,GHz, 7.125--8.4\,GHz and 14.8--15.35\,GHz.

\mdcfig{ch25_timeline}{The road to 6G as planned in late 2026. ITU-R Working Party 5D defines the
framework and evaluates candidate technologies; 3GPP, whose specification will almost certainly be the
main candidate, studies 6G in Release~20 and writes the first specifications in Release~21. Dates
after 2026 are plans, and plans in standardisation slip; read the right-hand half of the chart as
approximate.}

The technical work happens in 3GPP. After 5G-Advanced in Releases~18 (completed 2024) and~19
(completed around the end of 2025), \textbf{Release~20} contains both further 5G-Advanced
features and the first \emph{6G study items}. 3GPP held a dedicated 6G workshop in March 2025 and
approved the radio-access study in mid-2025; the service-requirements study of SA1 (on 6G use cases,
TR~22.870) had started earlier. A study item produces a technical report, not a specification: it
evaluates options and records agreements. The normative \textbf{Release~21} work, the first set of
6G specifications, is planned to start around mid-2027, with completion targeted around the end of
2028 or early 2029 and the first commercial deployments expected around 2030. As of September 2026
the study is in progress; the statements in this chapter about what 6G ``will'' contain are therefore
informed expectations, and are worded as such.

\begin{inpractice}[What the 6G study looks like from the inside]
A 3GPP study proceeds by agreements recorded in meeting reports, often worded as working
assumptions to be confirmed later. A few early directions were visible by 2026. CP-OFDM remains the
baseline downlink waveform and DFT-s-OFDM a strong uplink candidate, with new waveforms evaluated
against them; the design targets the upper mid-band as the main new spectrum while operating in
existing FR1 and FR2 bands; the system should share spectrum dynamically with 5G NR (multi-RAT
spectrum sharing) so that operators can introduce 6G in their existing bands; AI/ML-based features are
expected from the first release rather than retrofitted; sensing and non-terrestrial access are in
scope. Expect some of these to be refined by the time you read this. The authoritative record is the
3GPP technical reports and the chairs' notes of the RAN plenary and working-group meetings, all
freely downloadable from \texttt{3gpp.org}.
\end{inpractice}

\begin{pitfall}[Research targets are not requirements]
A paper that says ``6G will deliver 1\,Tb/s with 0.1\,ms latency'' is quoting a research aspiration,
sometimes a decade old. The values in M.2160 are explicitly \emph{examples}; the binding requirements
for IMT-2030 are those of the WP~5D technical-performance report, and even they are minimum
conditions for a candidate technology to be accepted, not promises about deployed networks. When you
cite a 6G number, cite the document and say which kind of number it is.
\end{pitfall}

%======================================================================
\section{Spectrum: the upper mid-band and the sub-terahertz frontier}
\label{sec:ch25:spectrum}
%======================================================================

\subsection{Why spectrum still matters}

From Chapter~\ref{ch:infotheory}, the capacity of a link is $B\log_2(1+\mathrm{SNR})$. At the SNRs of
a well-designed cellular network, doubling the bandwidth doubles the capacity, whereas doubling the
SNR adds about one bit per second per hertz. Every generation has therefore been accompanied by new
spectrum: 800--900\,MHz for 1G, 1.8--1.9\,GHz for 2G, 2.1\,GHz for 3G, 2.6\,GHz and the digital
dividend for 4G, and 3.3--4.2\,GHz plus 24--71\,GHz (FR2) for 5G. Table~\ref{tab:ch25:bands} lists the
ranges discussed for 6G.

\begin{table}[htbp]\centering\small
\caption{Frequency ranges in the 6G discussion. Bandwidths are indicative of what a single operator
might use; actual allocations depend on national regulators and the outcome of WRC-27.}
\label{tab:ch25:bands}
\begin{tabularx}{\linewidth}{@{}>{\raggedright\arraybackslash}p{2.2cm} >{\raggedright\arraybackslash}p{3.0cm} >{\raggedright\arraybackslash}p{1.7cm} >{\raggedright\arraybackslash}X@{}}
\toprule
Range & Example bands & Channel width & Role and main issues\\
\midrule
Sub-7\,GHz (FR1) & existing 4G/5G bands & 5--100\,MHz & Coverage layer; 6G introduced by spectrum sharing with NR\\
Upper mid-band (``FR3'') & 7.125--8.4, 12.7--13.25, 14.8--15.35\,GHz & 200--400\,MHz & Capacity layer on macro sites with large arrays; incumbents (satellite, radar, fixed links, science)\\
mmWave (FR2) & 24.25--71\,GHz & 0.4--2\,GHz & Hotspots, fixed wireless, backhaul; limited mobile deployment in 5G\\
D-band & 110--170\,GHz & 2--20\,GHz & Backhaul and fronthaul, short-range access, sensing\\
Low THz & 252--325\,GHz (IEEE 802.15.3d) & up to 69\,GHz & Kiosk downloads, chip-to-chip, data-centre links\\
\bottomrule
\end{tabularx}
\end{table}

\subsection{The upper mid-band}

The band from roughly 7 to 24\,GHz, sometimes called \term{FR3} after 5G's FR1 and FR2, has become the
centre of gravity of the 6G spectrum discussion. Its appeal is easiest to see through a worked
example.

\begin{worked}[Does a 15\,GHz macro cell cover like a 3.5\,GHz one?]
Take a 5G macro panel at 3.5\,GHz with 64 elements at half-wavelength spacing, about 35\,cm
square. Move to 15\,GHz with the \emph{same panel size}. The wavelength shrinks from 8.6\,cm to 2.0\,cm, so the number
of half-wavelength elements that fit in the same area grows by $(15/3.5)^2\approx18$, from 64 to about
1200, and the array gain grows by $10\log_{10}18\approx12.6$\,dB. The free-space path loss between
isotropic antennas rises by $20\log_{10}(15/3.5)=12.6$\,dB. The two cancel exactly: Friis's formula
in its aperture form,
\begin{equation}
\frac{P_r}{P_t}=\frac{A_tA_r}{\lambda^2d^2}\quad\text{(two apertures)},\qquad
\frac{P_r}{P_t}=\frac{G_rA_t}{4\pi d^2}\quad\text{(aperture to fixed-gain antenna)},
\label{eq:ch25:friisap}
\end{equation}
shows that with a fixed transmit aperture and a fixed-\emph{gain} (say, omnidirectional) receive
antenna, the received power does not depend on frequency at all. The handset's antenna, however, is a
fixed-gain antenna whose effective area shrinks as $\lambda^2$, and on the uplink the handset is the
transmitter. Diffraction around buildings, penetration through walls (several decibels more at
15\,GHz than at 3.5\,GHz for modern glass and concrete) and foliage loss also grow with frequency.
The upshot is that a 15\,GHz layer on existing macro sites can approach 3.5\,GHz downlink coverage
outdoors, needs an uplink helped by lower bands, and offers several hundred megahertz per operator.
The cost is in the panel: twelve hundred elements, with digital or hybrid beamforming, at macro power.
\end{worked}

The band is not empty. Parts of 7--8\,GHz are used by government satellites and fixed links,
12.7--13.25\,GHz by satellite feeder uplinks, around 10\,GHz and 15\,GHz by radars, and several
narrow bands by radio astronomy and Earth-exploration satellites whose receivers are extremely
sensitive. The 6G question in FR3 is therefore as much regulatory as technical: how to share with
incumbents (exclusion zones, dynamic access, coordination) and how quickly bands can be cleared.

\subsection{The sub-terahertz bands and molecular absorption}

Above 100\,GHz the spectrum is wide open: tens of gigahertz of contiguous bandwidth, enough for
100\,Gb/s with simple modulation. The physics is less friendly. Molecules in the atmosphere have
rotational energy levels whose transitions fall in the millimetre and sub-millimetre range, and a
wave at a transition frequency gives up energy to the gas. The oxygen molecule has a cluster of lines
around 60\,GHz and a single line at 119\,GHz; water vapour has lines at 22\,GHz, 183\,GHz, 325\,GHz and
many more above. Between the lines lie \emph{windows}: around 35, 94, 140, 220 and 300--340\,GHz.
The attenuation follows the Beer--Lambert law: power falls as $e^{-k(f)d}$, linear in decibels with
distance, at the \term{specific attenuation} $\gamma(f)$ in dB/km plotted in
Figure~\ref{fig:ch25_absorption}a. The authoritative model is ITU-R P.676, which sums the contributions
of dozens of spectral lines with pressure- and temperature-dependent widths.

\mdcfig{ch25_absorption}{(a) Approximate clear-air specific attenuation at sea level with
7.5\,g/m$^3$ of water vapour, from simplified formulas of earlier editions of ITU-R P.676 (valid to
about 350\,GHz). Shaded: the upper mid-band (FR3), FR2, the D-band and the IEEE 802.15.3d band. (b)
Path loss including antenna gains at 10\,m and 100\,m: between isotropic antennas (red) loss grows as
$f^2$, but between two fixed apertures of 5\,cm$\times$5\,cm (navy) it \emph{falls} as $f^2$ until the
absorption lines bite. At 100\,m the 183\,GHz and 325\,GHz lines are visible but modest; at 1\,km they
would dominate.}

Figure~\ref{fig:ch25_absorption}b makes the central point about high frequencies. Absorption at
140\,GHz (about 1\,dB/km) or 300\,GHz (about 5\,dB/km) is negligible over the 10--100\,m distances
for which these bands are proposed. What hurts is the isotropic free-space loss, $20\log_{10}(4\pi
d/\lambda)$, which grows by 20\,dB per decade of frequency---and that is entirely a statement about the
shrinking effective area of a fixed-gain antenna. Keep the apertures fixed and the trend reverses.
The sub-THz problem is thus not propagation loss as such but \emph{directivity}: an aperture of a few
centimetres at 300\,GHz has a gain of 40\,dBi and a beamwidth of about a degree, which must be pointed,
tracked through motion and re-acquired after blockage by a hand, a head or a passing person (body
blockage is typically 20--40\,dB at these frequencies). Rain adds about 10--20\,dB/km in heavy rain
(Chapter~\ref{ch:channels}), which matters for backhaul hops but not for a 10\,m link indoors.

\begin{worked}[A 300\,GHz link budget over 10\,m]
A kiosk download link at 300\,GHz uses one 51.84\,GHz channel of the IEEE 802.15.3d band; call it
50\,GHz.
Transmit power is 0\,dBm (realistic for an integrated 300\,GHz power amplifier, see below), and each
end has a lens or array antenna of 30\,dBi (an aperture of only about 1\,cm square, beamwidth about
6$^\circ$).
\begin{itemize}
\item Free-space loss: $\lambda=1.0$\,mm, $20\log_{10}(4\pi\times10/0.001)=102.0$\,dB.
\item Gaseous absorption: about 5\,dB/km$\times$0.01\,km $=0.05$\,dB. Negligible.
\item Pointing, polarisation and implementation losses: allow 3\,dB.
\item Received power: $0+30+30-102.0-0.05-3=-45.1$\,dBm.
\item Noise: $-174+10\log_{10}(5\times10^{10})=-67.0$\,dBm, plus a 10\,dB receiver noise figure
$=-57.0$\,dBm.
\item SNR $=-45.1-(-57.0)\approx12$\,dB.
\end{itemize}
The Shannon capacity is $50\times10^9\log_2(1+15.8)\approx204$\,Gb/s; with a practical 3\,dB gap to
capacity and a code rate around 0.8, about 100\,Gb/s is achievable with 16-QAM at a symbol rate near
40\,GBd. Change the distance to 100\,m and the SNR drops by 20\,dB (plus 0.5\,dB of absorption) to
$-8$\,dB, which still supports a few tens of gigabits per second at low spectral efficiency; but a
6$^\circ$ beam at 100\,m is 10\,m wide on one axis and only a few metres in the vertical if the
antenna is a slit, and a person walking through it costs 20\,dB or more. The link budget closes. The
system problem---beam tracking, blockage, and the power consumed by a 50\,GHz-wide receiver---is what
remains.
\end{worked}

\begin{history}[From Hertz to the terahertz gap]
Heinrich Hertz's experiments of 1887--88 used wavelengths of tens of centimetres to a few metres, and
Jagadish Chandra Bose in Calcutta demonstrated millimetre-wave apparatus---horn antennas, waveguides,
dielectric lenses and a galena crystal detector---around 1895--97, working at frequencies of tens of
gigahertz. Radio then went \emph{down} in frequency for decades, because long waves travelled further
and vacuum tubes could not generate microwaves efficiently. Radar in the Second World War opened the
centimetre bands, and the region between about 100\,GHz and 10\,THz became known as the ``terahertz
gap'', too high for electronics and too low for photonics. Silicon-germanium and indium-phosphide
transistors with maximum oscillation frequencies above 500\,GHz, and CMOS above 300\,GHz, have been
closing the gap since the 2010s. In 2017 IEEE 802.15.3d became the first wireless standard for the
252--325\,GHz range, and in 2019 the US FCC opened 95\,GHz--3\,THz for experimental licences and
designated over 20\,GHz of unlicensed spectrum above 95\,GHz.
\end{history}

\subsection{Hardware limits: power, noise and conversion}

Above about 100\,GHz the radio itself, rather than the channel, sets the limits. Three numbers
matter.

\paragraph{Transmit power and efficiency.} The maximum output power and efficiency of a transistor
amplifier fall as the operating frequency approaches the device's $f_{\max}$. Typical integrated power
amplifiers deliver tens of watts with 40--50\% efficiency at 3.5\,GHz (in GaN), about 20--25\,dBm per
element at 28\,GHz, roughly 10--20\,dBm at D-band and around 0--10\,dBm at 300\,GHz, with
power-added efficiencies falling to the low tens of per cent and below. Large arrays compensate by
summing many small amplifiers in space---the EIRP of an $N$-element array grows as $N^2$ ($N$ times the
power and $N$ times the gain)---but the DC power grows as $N$, and the heat must be removed from a
small area.

\paragraph{Phase noise.} A local oscillator at $f_c$ is usually synthesised by multiplying a
lower-frequency reference, and frequency multiplication by $K$ multiplies phase deviations by $K$,
raising the phase-noise spectrum by $20\log_{10}K$\,dB. A synthesiser that achieves
$-100$\,dBc/Hz at 1\,MHz offset at 28\,GHz would, built the same way, give about $-80$\,dBc/Hz at
280\,GHz. Chapter~\ref{ch:ofdm} showed that phase noise causes a common phase error, which pilots
remove, and inter-carrier interference, which they do not. The ICI is the part of the phase-noise
spectrum that lies beyond about one subcarrier spacing, so the remedy is wider subcarriers, as in
Figure~\ref{fig:ch25_hardware}b. That figure contains an important subtlety: widening the subcarrier
spacing helps only once it exceeds the bandwidth of the synthesiser's phase-locked loop, inside which
the phase noise is flat. This is one reason why NR introduced phase-tracking reference signals and
120--960\,kHz numerologies for FR2, and why sub-THz proposals use subcarrier spacings of several
megahertz or single-carrier waveforms.

\paragraph{Data conversion.} The power of an analog-to-digital converter is summarised by the
\term{Walden figure of merit}
\begin{equation}
\mathrm{FOM}_W=\frac{P}{2^{\mathrm{ENOB}}f_s}\quad\text{(joules per conversion step)},
\label{eq:ch25:walden}
\end{equation}
where ENOB is the effective number of bits and $f_s$ the sample rate. The best converters in
B.~Murmann's long-running survey reach a few femtojoules per step at moderate rates; at tens of
gigasamples per second, tens of femtojoules is good. Equation~\eqref{eq:ch25:walden} says that power
doubles with every extra bit and with every doubling of bandwidth (Figure~\ref{fig:ch25_hardware}a).

\mdcfig{ch25_hardware}{Two hardware limits of wideband 6G radios. (a) ADC power from the Walden
figure of merit, assuming a constant 50\,fJ per conversion step (real converters degrade above a
few GS/s, so this is optimistic at the right of the plot). (b) The SNR ceiling set by phase-noise ICI
after common-phase-error removal, for a toy synthesiser with $-100$\,dBc/Hz at 1\,MHz offset at
28\,GHz, a 500\,kHz loop bandwidth and a white floor, scaled by $20\log_{10}(f_c/28\,\text{GHz})$.
Wider subcarriers help only beyond the loop bandwidth.}

\begin{worked}[The power of a fully digital sub-THz array]
A 140\,GHz access point has 256 antenna elements, each with its own pair of ADCs (I and Q) sampling at
10\,GS/s with 6 effective bits, and $\mathrm{FOM}_W=50$\,fJ/step. Each ADC dissipates
$50\times10^{-15}\times2^6\times10^{10}=32$\,mW, and the array's 512 converters dissipate 16.4\,W---before
any digital signal processing, which at 10\,GS/s per element and a few operations per sample is far
larger still. Halving the resolution to 3 bits saves a factor of 8 in converter power; this is why
low-resolution (even one-bit) converters with massive arrays are a serious research topic, and why
practical sub-THz radios use hybrid beamforming (Chapter~\ref{ch:mimo}) with only a few digital
chains. The same arithmetic, at 400\,MHz and 12 bits, explains why a 5G massive-MIMO radio can afford
64 digital chains.
\end{worked}

\begin{keyidea}[The radio becomes the bottleneck]
At sub-6\,GHz the channel is the scarce resource and signal processing is cheap. At sub-THz the channel
is abundant and the transistors are the scarce resource: output power, phase noise and converter
energy, all of which worsen with frequency and bandwidth. Waveforms for these bands are chosen for
the hardware---low peak-to-average ratio, tolerance of phase noise, low-resolution conversion---rather
than for spectral efficiency.
\end{keyidea}

%======================================================================
\section{New waveforms for a fast-moving world}
\label{sec:ch25:waveforms}
%======================================================================

OFDM has carried every major wireless standard of the last twenty years (Chapter~\ref{ch:ofdm}), and
it will almost certainly be the baseline of 6G. Its virtues are hard to beat: one-tap equalisation of
frequency-selective channels, trivially flexible resource allocation, straightforward MIMO, and a
mature ecosystem. Its weaknesses are also familiar: a high peak-to-average power ratio, sensitivity to
frequency offset and phase noise, and the assumption that the channel is constant over each symbol.
This section examines the waveforms proposed to address the last of these, the regime of very high
Doppler, and then looks briefly at other ideas.

\subsection{The doubly dispersive channel revisited}

Chapter~\ref{ch:channels} described a time-varying multipath channel by its \emph{spreading
function} $h(\tau,\nu)$, the complex gain of the scatterers at delay $\tau$ and Doppler $\nu$:
\begin{equation}
r(t)=\iint h(\tau,\nu)\,s(t-\tau)\,e^{j2\pi\nu(t-\tau)}\,d\tau\,d\nu+w(t).
\label{eq:ch25:spreading}
\end{equation}
For $P$ discrete paths, $h(\tau,\nu)=\sum_ph_p\delta(\tau-\tau_p)\delta(\nu-\nu_p)$. A path's delay
and Doppler are set by geometry---distance and the rate of change of distance---and they change
slowly, over the time it takes a vehicle to move a significant fraction of the distance to the
scatterer. The time--frequency response $H(f,t)$ that OFDM sees, in contrast, is the two-dimensional
Fourier transform of $h(\tau,\nu)$ and fluctuates on the scale of the coherence time $\sim1/\nu_{\max}$
and the coherence bandwidth $\sim1/\tau_{\max}$ (Figure~\ref{fig:ch25_dd_channel}a).

OFDM with subcarrier spacing $\Delta f$ and symbol duration $T=1/\Delta f$ needs $\tau_{\max}$ to fit in
the cyclic prefix and $\nu_{\max}T\ll1$. When the second condition fails, the Doppler spread smears each
subcarrier into its neighbours. For a classical (Jakes) Doppler spectrum the inter-carrier interference
power is bounded by
\begin{equation}
P_{\text{ICI}}\lesssim\frac{(\pi\nu_{\max}T)^2}{3},
\label{eq:ch25:ici}
\end{equation}
so the signal-to-ICI ratio falls by 6\,dB for every doubling of speed. Wider subcarriers reduce $T$,
but then a fixed cyclic prefix becomes a larger fraction of the symbol. High-speed trains (up to
500\,km/h), aircraft (up to 1000\,km/h) and low-orbit satellites (Chapter~\ref{ch:satellite}) push
OFDM to the edge of this trade.

\subsection{OTFS: modulating in the delay--Doppler domain}

\term{Orthogonal time frequency space} (OTFS) modulation, introduced by Hadani and colleagues in 2017,
places the data symbols on a grid in the delay--Doppler plane instead of the time--frequency plane.
Take an OFDM-like frame of $N$ symbols of duration $T$, each with $M$ subcarriers spaced
$\Delta f=1/T$; the frame lasts $NT$ and occupies $M\Delta f$. The reciprocal grid in delay--Doppler has
$M$ delay bins of width $1/(M\Delta f)$ and $N$ Doppler bins of width $1/(NT)$. The data
$x[\ell,k]$, $\ell=0,\dots,M-1$ (delay), $k=0,\dots,N-1$ (Doppler), are mapped to the time--frequency grid
by the \term{inverse symplectic finite Fourier transform} (ISFFT),
\begin{equation}
X[m,n]=\frac{1}{\sqrt{MN}}\sum_{k=0}^{N-1}\sum_{\ell=0}^{M-1}x[\ell,k]\,
e^{\,j2\pi\left(\frac{nk}{N}-\frac{m\ell}{M}\right)},
\label{eq:ch25:isfft}
\end{equation}
an inverse DFT along Doppler (to time index $n$) and a DFT along delay (to frequency index $m$). The
time--frequency samples are then transmitted with a multicarrier modulator, the \term{Heisenberg
transform}
\begin{equation}
s(t)=\sum_{n=0}^{N-1}\sum_{m=0}^{M-1}X[m,n]\,g_{\text{tx}}(t-nT)\,e^{\,j2\pi m\Delta f(t-nT)},
\label{eq:ch25:heisenberg}
\end{equation}
which for a rectangular pulse $g_{\text{tx}}$ is exactly OFDM (without per-symbol prefixes; OTFS
usually uses one prefix for the whole frame). The receiver reverses both steps: a bank of matched
filters (the \emph{Wigner transform}, which for rectangular pulses is the OFDM DFT) produces
$Y[m,n]$, and the \term{symplectic finite Fourier transform} (SFFT), the inverse of
\eqref{eq:ch25:isfft}, returns to the delay--Doppler domain. Figure~\ref{fig:ch25_otfs_chain} shows the
chain.

\begin{figure}[htbp]\centering
\begin{tikzpicture}[node distance=0.38cm and 0.38cm,
  blk/.style={draw=navy,fill=navy!6,thick,minimum height=0.95cm,minimum width=1.4cm,align=center,font=\sffamily\scriptsize},
  dom/.style={font=\sffamily\tiny,text=accent,align=center},
  arr/.style={-{Stealth[length=2mm]},thick,navy}]
\node[blk] (x) {DD symbols\\$x[\ell,k]$};
\node[blk,right=of x] (isfft) {ISFFT\\(DFT on $\ell$,\\IDFT on $k$)};
\node[blk,right=of isfft] (heis) {Heisenberg\\(OFDM mod.)\\$+$ one CP};
\node[blk,right=of heis,fill=accent!6,draw=accent] (ch) {channel\\$h(\tau,\nu)$};
\node[blk,right=of ch] (wig) {Wigner\\(OFDM\\demod.)};
\node[blk,right=of wig] (sfft) {SFFT};
\node[blk,right=of sfft] (det) {DD\\detector};
\draw[arr] (x)--(isfft); \draw[arr] (isfft)--(heis); \draw[arr] (heis)--(ch); \draw[arr] (ch)--(wig); \draw[arr] (wig)--(sfft); \draw[arr] (sfft)--(det);
\node[dom,above=0.12cm of isfft] {DD $\to$ TF}; \node[dom,above=0.12cm of heis] {TF $\to$ time};
\node[dom,above=0.12cm of wig] {time $\to$ TF}; \node[dom,above=0.12cm of sfft] {TF $\to$ DD};
\draw[decorate,decoration={brace,amplitude=5pt,mirror},navy!70] ([yshift=-3pt]isfft.south west) -- ([yshift=-3pt]heis.south east)
  node[midway,below=6pt,font=\sffamily\scriptsize,text=navy,align=center]{rectangular pulses: $\vect{s}=\operatorname{vec}(\vect{x}\vect{F}_N^{\mathsf H})$\\(an $N$-point IDFT per delay bin)};
\draw[decorate,decoration={brace,amplitude=5pt,mirror},navy!70] ([yshift=-3pt]wig.south west) -- ([yshift=-3pt]sfft.south east)
  node[midway,below=6pt,font=\sffamily\scriptsize,text=navy]{an existing OFDM receiver + a 2-D FFT};
\end{tikzpicture}
\caption{The OTFS transceiver as a two-dimensional precoder around an OFDM modem. With rectangular
pulses the ISFFT and the OFDM modulator collapse into one $N$-point inverse DFT along the Doppler
axis, and the transmitted frame is simply the columns of $\vect{x}\vect{F}_N^{\mathsf H}$ sent one
after another---the discrete Zak-transform view of OTFS.}
\label{fig:ch25_otfs_chain}
\end{figure}

\subsubsection{The matrix picture}
The cleanest way to see what OTFS does is with matrices, using the tools of Chapter~\ref{ch:ofdm}.
Arrange the DD symbols as an $M\times N$ matrix $\vect{x}$. The ISFFT is
$\vect{X}=\vect{F}_M\vect{x}\vect{F}_N^{\mathsf H}$, with $\vect{F}_K$ the unitary $K$-point DFT
matrix, and OFDM modulation of each column is multiplication by $\vect{F}_M^{\mathsf H}$. The
transmitted block is therefore
\begin{equation}
\vect{S}=\vect{F}_M^{\mathsf H}\vect{F}_M\vect{x}\vect{F}_N^{\mathsf H}=\vect{x}\vect{F}_N^{\mathsf H},
\qquad
\vect{s}=\operatorname{vec}(\vect{S})=(\vect{F}_N^{\mathsf H}\otimes\vect{I}_M)\,\operatorname{vec}(\vect{x}),
\label{eq:ch25:otfsvec}
\end{equation}
read out column by column, with one cyclic prefix of at least $\ell_{\max}$ samples in front of the
$MN$-sample frame. The two DFTs along the delay axis cancel: OTFS is an $N$-point IDFT along Doppler
for each delay bin, which spreads each symbol over the whole frame in time.

Now pass $\vect{s}$ through a channel with $P$ paths, path $p$ having integer delay $\ell_p$ samples and
Doppler $\nu_p$ in cycles per sample. With sampling interval $1/(M\Delta f)$, a Doppler of one DD bin,
$1/(NT)$\,Hz, is $1/(MN)$ cycles per sample. The time-domain channel acting on the prefixed frame is
\begin{equation}
\vect{H}=\sum_{p=1}^{P}h_p\,\bm{\Delta}^{k_p}\bm{\Pi}^{\ell_p},\qquad
\bm{\Pi}=\text{cyclic shift},\quad
\bm{\Delta}=\operatorname{diag}\big(e^{j2\pi n/(MN)}\big)_{n=0}^{MN-1},
\label{eq:ch25:hdd}
\end{equation}
where $k_p=MN\nu_p$ is the Doppler in DD bins. The effective channel seen by the DD symbols is
$\vect{H}_{\text{eff}}=(\vect{F}_N\otimes\vect{I}_M)\vect{H}(\vect{F}_N^{\mathsf H}\otimes\vect{I}_M)$.
Each term $\bm{\Delta}^{k}\bm{\Pi}^{\ell}$ becomes, after the transform, a two-dimensional cyclic
shift by $\ell$ in delay and $k$ in Doppler, multiplied by a known phase. When the delays and Dopplers
fall on the grid,
\begin{equation}
y[\ell,k]=\sum_{p=1}^{P}h_p\,\phi_p[\ell,k]\;x\big[[\ell-\ell_p]_M,\,[k-k_p]_N\big]+w[\ell,k],
\label{eq:ch25:ddio}
\end{equation}
with $|\phi_p|=1$: a \emph{two-dimensional circular convolution} of the data with the sparse channel.
Figure~\ref{fig:ch25_dd_channel}b shows it: a single symbol at the origin comes back as $P$ spikes at
the paths' delays and Dopplers. A path whose Doppler is not an integer number of bins leaks into
neighbouring Doppler bins with a Dirichlet-kernel profile, exactly as a tone between DFT bins does
(Chapter~\ref{ch:signals}); windowing and finer Doppler resolution (larger $N$) reduce the leakage.

\mdcfig{ch25_dd_channel}{The same three-path channel (delays 0, 3 and 7 samples; Dopplers 0, 3 and
$-4.6$ Doppler bins) in two domains, for $M=64$ subcarriers and $N=32$ symbols. (a) The time--frequency
response that OFDM must track varies on every symbol and every few subcarriers. (b) The response of
the DD domain to a single OTFS symbol: three spikes at the paths' coordinates; the path with a
fractional Doppler spreads along the Doppler axis.}

\subsubsection{Why OTFS tolerates high Doppler}
Three properties follow from~\eqref{eq:ch25:ddio}.
\begin{enumerate}
\item \textbf{Every symbol sees the whole channel.} Each DD symbol is spread over all $MN$
time--frequency resources, so it experiences every fade in the frame. A joint detector can collect
the energy of all $P$ resolvable paths, a diversity gain that OFDM obtains only through coding and
interleaving across the frame.
\item \textbf{The channel description is compact and stable.} The receiver needs only the $P$ triples
$(h_p,\ell_p,k_p)$, which change slowly. A single pilot in the DD grid, surrounded by a guard region
of empty bins, reveals all of them at once (the \emph{embedded pilot} of Raviteja, Phan and Hong,
2019). An OFDM receiver must instead estimate $MN$ time--frequency coefficients that change from
symbol to symbol.
\item \textbf{Doppler is not interference.} In OFDM, energy that leaks to the neighbouring subcarrier
is interference to be suppressed; in OTFS it is a known shift in a convolution, to be exploited.
\end{enumerate}
The price is a receiver that must undo a two-dimensional convolution: a linear MMSE equaliser on an
$MN\times MN$ sparse matrix, or iterative detectors such as message passing (Raviteja, Phan, Hong and
Viterbo, 2018), whose complexity grows with $MNP$. Latency is set by the frame length $NT$, since
nothing can be detected until the frame is complete.

\mdcfig{ch25_otfs_ber}{Uncoded QPSK on a four-path doubly dispersive channel (Rayleigh path gains,
delays up to 3 samples, Jakes-distributed Dopplers) with a $16\times16$ grid. (a) BER against SNR
at a maximum Doppler of 0.3 subcarrier spacings. (b) BER at 20\,dB against Doppler. The one-tap OFDM
equaliser has an ICI floor; an ICI-aware LMMSE equaliser for each OFDM symbol removes the floor but
has no diversity across the frame; OTFS and AFDM with LMMSE detection over the whole frame harvest
the path diversity. Generated with \texttt{commlib.sixg}.}

Figure~\ref{fig:ch25_otfs_ber} compares OTFS with OFDM by simulation. The one-tap OFDM equaliser,
which uses the average channel of each symbol, hits an error floor once the Doppler is a noticeable
fraction of the subcarrier spacing, as~\eqref{eq:ch25:ici} predicts. A per-symbol LMMSE equaliser
that knows the full ICI matrix removes the floor, but each OFDM symbol is still a single flat-fading
observation, so its error rate falls only as $1/\text{SNR}$. OTFS, detected jointly over the frame,
falls much faster.

\begin{pitfall}[An unfair race]
Comparisons like Figure~\ref{fig:ch25_otfs_ber} are common in the OTFS literature and easy to
misread. The OFDM curves are \emph{uncoded}, so they show no frequency or time diversity; the OTFS
receiver solves a $256\times256$ system per frame while the OFDM receiver solves sixteen $16\times16$
systems. Give the OFDM system a rate-1/2 code interleaved across the frame and an iterative receiver
and much of the gap closes; OTFS's advantage then lies mainly in channel estimation overhead and in
very high Doppler. Always ask what complexity, coding and channel knowledge each curve assumes.
\end{pitfall}

\begin{worked}[Sizing an OTFS grid for a 500\,km/h train at 4\,GHz]
A high-speed train at $v=500$\,km/h $=138.9$\,m/s, served at $f_c=4$\,GHz, sees a maximum Doppler
$\nu_{\max}=vf_c/c=138.9\times4\times10^9/3\times10^8=1852$\,Hz. Take a delay spread of $5\,\mu$s
(open terrain along the track).

\emph{OFDM baseline.} With $\Delta f=30$\,kHz, $T=33.3\,\mu$s and $\nu_{\max}T=0.062$. The ICI bound
\eqref{eq:ch25:ici} gives $(\pi\times0.062)^2/3=0.0126$, a signal-to-ICI ratio of about 19\,dB: adequate
for 16-QAM, marginal for 64-QAM, and worse if the train passes a trackside mast, where the Doppler
flips sign from $+\nu_{\max}$ to $-\nu_{\max}$ within a fraction of a second.

\emph{OTFS grid.} Keep $\Delta f=30$\,kHz and take $M=512$ subcarriers (15.36\,MHz). The delay
resolution is $1/(M\Delta f)=65$\,ns, so the 5\,$\mu$s delay spread spans $\ell_{\max}\approx77$ delay
bins. With $N=64$ symbols the frame lasts $NT=2.13$\,ms and the Doppler resolution is
$1/(NT)=469$\,Hz, so $\nu_{\max}$ spans $k_{\max}\approx4$ Doppler bins on each side. Both are well
inside the grid ($77\ll512$, $2\times4+1\ll64$), so the channel is resolvable and sparse.

\emph{Pilot overhead.} An embedded pilot needs a guard region extending $\ell_{\max}$ bins either side
in delay and $2k_{\max}$ either side in Doppler, about $(2\times77+1)(4\times4+1)=2635$ bins, or 8\% of the
32\,768-bin frame, for a channel estimate valid over the whole frame regardless of speed.

\emph{Latency.} The frame takes 2.13\,ms; halving $N$ halves the latency but doubles the Doppler bin
width, so the Doppler spans only two bins and fractional-Doppler leakage becomes relatively worse.
Grid design is a trade between latency, resolution and pilot overhead.
\end{worked}

\subsection{AFDM: chirps instead of tones}

\term{Affine frequency division multiplexing} (AFDM), proposed by Bemani, Ksairi and Kountouris
(2023), takes a different route to the same goal. It replaces the DFT of OFDM by the \emph{discrete
affine Fourier transform} (DAFT),
\begin{equation}
\vect{A}=\bm{\Lambda}_{c_2}\vect{F}_N\bm{\Lambda}_{c_1},\qquad
\bm{\Lambda}_c=\operatorname{diag}\big(e^{-j2\pi cn^2}\big)_{n=0}^{N-1},
\label{eq:ch25:daft}
\end{equation}
and transmits $\vect{s}=\vect{A}^{\mathsf H}\vect{x}$ with a \emph{chirp-periodic} prefix (which reduces
to an ordinary cyclic prefix when $2Nc_1$ is an integer and $N$ is even). Each ``subcarrier'' is a
discrete chirp whose instantaneous frequency sweeps the whole band, wrapping around
(Figure~\ref{fig:ch25_afdm}a). With $c_1=c_2=0$, AFDM is OFDM.

The point of the chirps is what they do to a doubly dispersive channel. A path with delay $\ell_p$ and
integer Doppler $\alpha_p$ (in units of $1/N$ cycles per sample) maps each DAFT-domain symbol to a
single other one, shifted by $(\alpha_p+2Nc_1\ell_p)\bmod N$. If the chirp rate is chosen as
\begin{equation}
c_1=\frac{2(\alpha_{\max}+\xi)+1}{2N},
\label{eq:ch25:c1}
\end{equation}
with $\xi$ a small guard for fractional Doppler, paths with different delays land in non-overlapping
ranges of the DAFT domain, and paths with the same delay but different Doppler are separated within
each range. Every path is resolved, and the effective channel $\vect{A}\vect{H}\vect{A}^{\mathsf H}$ is a
set of separated diagonals (Figure~\ref{fig:ch25_afdm}c), the one-dimensional analogue of OTFS's
two-dimensional sparsity. Bemani and colleagues showed that AFDM achieves full diversity in linear
time-varying channels when, roughly, $(\ell_{\max}+1)\big(2(\alpha_{\max}+\xi)+1\big)\le N$, so that the
paths' ranges fit side by side in the DAFT domain (with appropriate $c_2$). The
second parameter, $c_2$, is a free phase rotation that can improve diversity or carry information.

\mdcfig{ch25_afdm}{AFDM in pictures, for $N=64$ and a three-path channel with delays 0, 2 and 5 and
Dopplers 0, 1 and $-2$ bins. (a) Instantaneous frequency of three AFDM chirp subcarriers. (b) The
OFDM effective channel $\vect{F}\vect{H}\vect{F}^{\mathsf H}$: Doppler spreads energy from the
diagonal into its neighbours (inter-carrier interference). (c) The AFDM effective channel
$\vect{A}\vect{H}\vect{A}^{\mathsf H}$: each path becomes its own, well-separated diagonal.}

AFDM is attractive for standards because it keeps the one-dimensional, single-DFT structure of
OFDM---an AFDM modem is an OFDM modem with two diagonal phase multiplications---and because its
chirps are a natural radar waveform (Section~\ref{sec:ch25:isac}). Its performance in
Figure~\ref{fig:ch25_otfs_ber} is very close to that of OTFS; the practical differences lie in pilot
design, MIMO integration, peak-to-average ratio and the cost of the receiver.

\begin{inpractice}[Will 6G use OTFS or AFDM?]
Probably not as its main waveform, at least not in the first release. OFDM's ecosystem advantage is
enormous, and OFDM itself can be made more Doppler-robust with wider subcarriers, more frequent
reference signals and better receivers; NR already supports 120\,kHz and wider numerologies and
tracking reference signals. Delay--Doppler processing is more likely to arrive in pieces: DD-domain
channel estimation and prediction inside an OFDM receiver, OTFS-like precoding as an option for
high-speed or satellite scenarios, and DD-domain processing for sensing. Because OTFS can be built
as a precoder on top of an OFDM modem (Figure~\ref{fig:ch25_otfs_chain}), such options can be added
later without a new physical layer. The same is true of AFDM.
\end{inpractice}

\subsection{Other waveform ideas}

\paragraph{Evolving DFT-s-OFDM.} The single-carrier-like DFT-spread OFDM of the NR uplink
(Chapter~\ref{ch:ofdm}) is the natural choice where the power amplifier dominates: coverage-limited
uplinks, satellites and sub-THz. Proposals to improve it include frequency-domain spectral shaping
(a mild RRC-like taper across the occupied subcarriers, which lowers the PAPR of $\pi/2$-BPSK and QPSK
further), zero-tail and unique-word variants that replace the cyclic prefix by low-power samples
generated inside the DFT, and constant-envelope schemes for the most power-limited links.

\paragraph{Faster-than-Nyquist revisited.} Chapter~\ref{ch:baseband} showed that binary sinc pulses
can be packed at $\tau\approx0.8$ of the Nyquist spacing without loss of minimum distance (the Mazo
limit). In multicarrier form, \emph{time--frequency packing} squeezes both the symbol period and the
subcarrier spacing below the orthogonal values. The gain---up to about 25\% more symbols in the same
bandwidth---is real but must be bought with a receiver that handles controlled ISI and ICI, and it
competes with the simpler alternative of a denser constellation and a stronger code. Its most
plausible home is where the amplifier forces a low-order constellation, such as satellite downlinks.

\paragraph{Index modulation.} In \term{index modulation} some information is carried by \emph{which}
resources are active. In OFDM with index modulation, each group of $n$ subcarriers activates $k$ of
them, carrying $\lfloor\log_2\binom{n}{k}\rfloor$ bits in the pattern plus $k\log_2M$ in the
constellation symbols; with $n=4$, $k=2$ and QPSK, a group carries $2+4=6$ bits on four subcarriers, fewer
than the 8 bits of plain QPSK, but with the energy concentrated on half the subcarriers, which helps at
low rates and in fast fading. \emph{Spatial modulation} (Mesleh and colleagues, 2008) activates one
of $N_t$ transmit antennas, carrying $\log_2N_t$ extra bits with a single RF chain. Both trade
spectral efficiency for energy efficiency or hardware simplicity, and both have so far remained
research topics; they are worth knowing as examples of a general idea---information can be encoded in
any degree of freedom the receiver can distinguish.

%======================================================================
\section{Extreme MIMO: from massive arrays to smart surfaces}
\label{sec:ch25:mimo}
%======================================================================

Chapter~\ref{ch:mimo} developed multi-antenna theory from diversity combining to massive MIMO. Two
trends push it further in 6G. Arrays are becoming physically \emph{large} compared with the distance
to the user, so the plane-wave model that underlies steering vectors and codebooks breaks down. And
antennas are being \emph{distributed}---over many access points, over walls and windows, over the
surfaces of buildings---so the notion of a single array, and of a cell, dissolves.

\subsection{Near field and the Rayleigh distance}

Place an $N$-element linear array of length $D$ along the $x$-axis, centred on the origin, and a user
at range $r$ and angle $\theta$ from broadside. The distance from element $n$ at position $x_n$ to the
user is, exactly,
\begin{equation}
d_n=\sqrt{r^2+x_n^2-2rx_n\sin\theta}\approx r-x_n\sin\theta+\frac{x_n^2\cos^2\theta}{2r},
\label{eq:ch25:fresnel}
\end{equation}
where the approximation (the \term{Fresnel approximation}) keeps the second-order term of the Taylor
expansion. The familiar steering vector of Chapter~\ref{ch:mimo} keeps only the linear term: the
wavefront is treated as a plane, and the phase across the array depends only on the angle. The
quadratic term is the curvature of a spherical wavefront, and it depends on the range. It can be
neglected when the phase it contributes at the edges of the array, $x_n=\pm D/2$, is small. At
broadside the phase error is $\pi D^2/(4\lambda r)$; requiring it to be at most $\pi/8$ gives the
classical \term{Rayleigh distance} (or Fraunhofer distance)
\begin{equation}
d_R=\frac{2D^2}{\lambda}.
\label{eq:ch25:rayleigh}
\end{equation}
Beyond $d_R$ an antenna is in its \emph{far field} (the Fraunhofer region): its pattern depends only on
direction. Between roughly $0.62\sqrt{D^3/\lambda}$ and $d_R$ lies the radiating \emph{near field}
(the Fresnel region), where the pattern depends on range as well. Closer still is the reactive near
field, which matters for antenna design but not for links.

\begin{worked}[The Rayleigh distance of a 1\,m array at 30\,GHz]
At 30\,GHz, $\lambda=1.0$\,cm. A 1\,m aperture---a 200-element half-wavelength linear array, or the
long side of a large panel---has $d_R=2\times1^2/0.01=200$\,m. Every user in a typical urban small cell
is in its near field. The same aperture at 3.5\,GHz ($\lambda=8.6$\,cm) has $d_R=23$\,m; at 140\,GHz,
933\,m. For a user at an angle $\theta$, the relevant aperture is the projected length $D\cos\theta$,
so the near-field region shrinks away from broadside. Figure~\ref{fig:ch25_rayleigh} generalises the
calculation: as frequencies rise and arrays grow, the near field moves from an antenna-measurement
curiosity to the normal operating regime.
\end{worked}

\mdcfig[0.72]{ch25_rayleigh}{Rayleigh distance $2D^2/\lambda$ against aperture size. The shaded band
marks typical small-cell and indoor link distances. Large arrays in the upper mid-band and beyond put
most users in the radiating near field.}

\subsection{Beamfocusing}

In the far field, a beamformer that matches the linear phase of a direction produces a \emph{beam}:
a wedge of high gain extending to infinity. In the near field, matching the exact spherical phases
$e^{\,j2\pi d_n/\lambda}$ of a point produces a \emph{spot}: the signals add coherently only near the
focal point, both in angle and in range (Figure~\ref{fig:ch25_nearfield}). This is
\term{beamfocusing}, the same principle as an optical lens. Its depth of focus can be estimated
from~\eqref{eq:ch25:fresnel}: for a focal range $r_F$ well inside the Rayleigh distance
($r_F\lesssim d_R/10$), the 3\,dB gain region extends approximately from
\begin{equation}
\frac{d_Rr_F}{d_R+10r_F}\quad\text{to}\quad\frac{d_Rr_F}{d_R-10r_F},
\label{eq:ch25:depth}
\end{equation}
an approximation due to Bj\"ornson, Demir and Sanguinetti. For the 1.37\,m array of
Figure~\ref{fig:ch25_nearfield} at 28\,GHz ($d_R=351$\,m) focused at 6\,m, it gives 5.1 to 7.2\,m;
the simulation in panel (c) shows 5.3 to 6.9\,m. As $r_F$ approaches $d_R/10$ the upper limit runs off
to infinity and the focus degenerates into an ordinary beam.

\mdcfig{ch25_nearfield}{Near-field beamfocusing with a 256-element half-wavelength array at 28\,GHz
(aperture 1.37\,m, white bar at the bottom of panels a and b; Rayleigh distance 351\,m). (a) A
conventional far-field beam steered to 20$^\circ$ is a wedge. (b) Weights matched to the spherical
wavefront of a point 6\,m away concentrate energy in a spot. (c) Gain along the 20$^\circ$ direction
for three focal distances and for the far-field beam: the focused beams have a finite depth that
grows with focal range; the far-field beam reaches full gain only near the Rayleigh distance.}

Beamfocusing has three consequences for system design. First, \emph{far-field codebooks lose gain}
in the near field: a beam steered at the right angle can be 10\,dB or more below a focused one at
short range (Figure~\ref{fig:ch25_nearfield}c), so codebooks and channel estimators must sample range
as well as angle---the ``polar-domain'' codebooks proposed for near-field estimation. Second, users at
the same angle but different ranges can be \emph{separated}, which multiplies the spatial degrees of
freedom available to multi-user MIMO: a new kind of spatial multiplexing in distance. Third,
energy can be delivered to a point with little spill-over, which matters for wireless power transfer
and for limiting exposure and interference.

\subsection{XL-MIMO and holographic MIMO}

\term{Extremely large-scale MIMO} (XL-MIMO) is the general name for arrays so large that, besides
near-field effects, the channel is \emph{spatially non-stationary}: different parts of the array see
different scatterers and different users, and a user may be visible from only part of the array.
Such arrays could be built into the facade of a building or along a stadium roof, with thousands of
elements. Processing then naturally becomes distributed along the array, with subarrays handling
their own visibility regions.

\term{Holographic MIMO} takes the limit of an array with a very large number of very closely spaced
elements, approaching a continuous aperture whose surface current can be shaped at will. The physics
that limits it is old. Pizzo, Marzetta and Sanguinetti showed that for a planar aperture of area $A$ in
isotropic scattering the number of spatial degrees of freedom is approximately
\begin{equation}
N_{\text{DoF}}\approx\frac{\pi A}{\lambda^2},
\label{eq:ch25:dofarea}
\end{equation}
which is reached already with half-wavelength spacing. Closer spacing adds no new degrees of freedom;
it can improve pattern control and, with mutual coupling exploited deliberately, give
\emph{superdirectivity} (gain above the usual $4\pi A/\lambda^2$), but superdirective arrays are
notoriously narrowband and lossy. In line of sight, the degrees of freedom between two parallel linear
apertures of lengths $L_t$ and $L_r$ at distance $d$ are approximately
\begin{equation}
N_{\text{DoF}}\approx\frac{L_tL_r}{\lambda d}+1,
\label{eq:ch25:doflos}
\end{equation}
the same physics as the line-of-sight MIMO design rule~\eqref{eq:ch19:losmimo} of
Chapter~\ref{ch:mimo}. Two 1\,m arrays at 30\,GHz, 10\,m apart, support about ten parallel streams in
pure line of sight; at 200\,m (the Rayleigh distance) only about one. The near field is where line-of-sight
MIMO lives.

\subsection{Cell-free massive MIMO}

Chapter~\ref{ch:mimo} introduced \term{cell-free massive MIMO}: many access points (APs), each with
a few antennas, distributed over an area and jointly serving all users through a central processing
unit (CPU). It addresses the oldest problem of cellular networks: the cell edge. In a cellular
network a user's performance depends on where it is relative to its serving site and the interfering
ones; a user near a cell boundary suffers both weak signal and strong interference. In a cell-free
network, every user is near \emph{some} APs, and the APs that would have been interferers become
helpers.

\subsubsection{System model and processing}
Let $L$ APs with $N$ antennas each serve $K$ single-antenna users. On the uplink, AP $l$ receives
$\vect{y}_l=\sum_{k=1}^{K}\vect{h}_{kl}s_k+\vect{n}_l$, where $\vect{h}_{kl}\in\C^N$ is the channel
from user $k$, with a large-scale gain $\beta_{kl}$ that varies over orders of magnitude across the
APs. Stack everything into $\vect{y}=\sum_k\vect{h}_ks_k+\vect{n}$ with $\vect{h}_k\in\C^{LN}$. If all
the received signals reach the CPU, it can apply the MMSE combiner of Chapter~\ref{ch:mimo} to the
whole network, giving user $k$ the SINR
\begin{equation}
\mathrm{SINR}_k=p_k\vect{h}_k^{\mathsf H}\Big(\sum_{i\ne k}p_i\vect{h}_i\vect{h}_i^{\mathsf
H}+\sigma^2\vect{I}_{LN}\Big)^{-1}\vect{h}_k.
\label{eq:ch25:cfsinr}
\end{equation}
This is \emph{centralised} operation. At the other extreme, each AP estimates the channels and
combines locally (for example by matched filtering, $\hat s_{kl}=\vect{h}_{kl}^{\mathsf H}\vect{y}_l$, as
in the original cell-free proposal of Ngo and colleagues, 2017), and the CPU merely adds the local
estimates, possibly with optimised weights. Bj\"ornson and Sanguinetti (2020) classified four levels
of cooperation between these extremes and showed that the centralised and the best distributed
schemes can differ by a factor of two or more in spectral efficiency, with the small-cell network
(each user served by one AP, no cooperation) far behind both.

\mdcfig{ch25_cellfree}{Cell-free massive MIMO against the alternatives with the same 64 antennas, in a
500\,m$\times$500\,m area with 16 users (3GPP-like urban path loss with 4\,dB shadowing, 20\,MHz,
100\,mW per user, perfect channel knowledge, uplink). (a) Layout: 64 single-antenna APs on a grid,
or one 64-antenna site in the centre. (b) Distribution of per-user spectral efficiency: the
co-located massive-MIMO site serves central users well and edge users poorly; uncoordinated small
cells suffer interference; centralised cell-free MMSE gives every user a strong signal and
suppresses interference.}

Figure~\ref{fig:ch25_cellfree} reproduces the qualitative result of the early cell-free papers. With
the same total number of antennas, the cell-free network raises the 5th-percentile spectral efficiency
(the cell-edge rate) by a large factor over both the co-located array and the small cells. Perfect
channel knowledge flatters every curve; with pilot contamination and estimation error the numbers
fall, but the ranking survives in the literature's more careful studies.

\subsubsection{Making it scale}
A network in which every AP serves every user does not scale: the CPU's computation and the
fronthaul traffic grow without bound as the network grows. Practical designs are
\emph{user-centric}: each user is served by a cluster of the APs with the strongest channels to it
(\emph{dynamic cooperation clusters}), and each AP serves a bounded number of users. The cluster
moves with the user, so there are no handovers in the cellular sense. The remaining engineering
problems are formidable:
\begin{itemize}
\item \textbf{Fronthaul.} Shipping raw IQ samples is expensive: an AP with four antennas on a
100\,MHz carrier at 122.88\,MS/s with 16-bit I and Q samples produces
$4\times122.88\times10^6\times32\approx15.7$\,Gb/s. Splitting the processing so that APs send
per-user soft estimates instead (the ``functional split'' question of Chapter~\ref{ch:cellular})
reduces this drastically, at some loss in interference suppression.
\item \textbf{Synchronisation.} Coherent joint transmission on the downlink requires the APs to be
phase-synchronised and their transmit and receive chains reciprocity-calibrated, across a network
rather than within one radio.
\item \textbf{Deployment.} Running fibre and power to hundreds of small APs costs more than one large
site. Ericsson's ``radio stripes'' concept---APs daisy-chained along a cable that carries power,
fronthaul and synchronisation, hidden in a building's fabric---is one answer.
\end{itemize}
Cell-free ideas are likely to enter 6G as enhanced multi-TRP (multiple transmission and reception
point) operation and distributed MIMO within operator-controlled clusters, rather than as a
city-wide network without cells.

\subsection{Reconfigurable intelligent surfaces}

A \term{reconfigurable intelligent surface} (RIS) is a panel of $N$ passive elements, each of which
reflects an incident wave with a controllable phase shift $\theta_n$ (Chapter~\ref{ch:mimo}). The
received signal with a direct path $h_d$ and the RIS path is
\begin{equation}
y=\Big(h_d+\sum_{n=1}^{N}g_n\,e^{j\theta_n}\,h_n\Big)x+w,
\label{eq:ch25:ris}
\end{equation}
where $h_n$ is the channel from the transmitter to element $n$ and $g_n$ from element $n$ to the
receiver. Choosing $\theta_n=\arg h_d-\arg(g_nh_n)$ aligns all the reflections with the direct path,
and the RIS term grows in amplitude as $N$ and in power as $N^2$: the same coherent array gain as a
phased array, obtained without any amplifiers or radio chains.

\subsubsection{The product-distance law}
How large is each term? Treat the optimally phased RIS in the far field as an aperture of area $NA$
(with $A$ the area of one element, typically $(\lambda/2)^2$). It \emph{captures} the power
$P_tG_tNA/(4\pi d_1^2)$ from a transmitter at distance $d_1$, and re-radiates it as an aperture with
gain $4\pi NA/\lambda^2$ towards a receiver at distance $d_2$ whose antenna has effective area
$G_r\lambda^2/4\pi$. Multiplying,
\begin{equation}
\frac{P_r}{P_t}\bigg|_{\text{RIS}}=\frac{G_tG_r\,(NA)^2}{(4\pi)^2\,d_1^2\,d_2^2}
\qquad\text{versus}\qquad
\frac{P_r}{P_t}\bigg|_{\text{direct}}=G_tG_r\Big(\frac{\lambda}{4\pi d}\Big)^2,
\label{eq:ch25:productlaw}
\end{equation}
ignoring the obliquity factors of the incident and reflected angles. The RIS path decays with the
\emph{product} of the squared distances, $d_1^2d_2^2$, instead of the square of the total distance:
the \term{product-distance law}, also called the double path loss. Setting the two equal gives the
aperture at which the RIS path matches a line-of-sight direct path:
\begin{equation}
NA=\frac{\lambda d_1d_2}{d}.
\label{eq:ch25:riseq}
\end{equation}
Equation~\eqref{eq:ch25:riseq} has two lessons (Figure~\ref{fig:ch25_ris}). The product $d_1d_2$ is
smallest when the RIS is close to one end of the link, so an RIS belongs near the transmitter or,
more usefully, near the users it serves. And with half-wavelength elements, $A=\lambda^2/4$, the
required \emph{number} of elements $N=4d_1d_2/(\lambda d)$ grows in proportion to frequency: an RIS
of a given physical size has more elements at higher frequency, but needs them.

The law changes when the RIS is large and close. If the user is within the RIS's own Rayleigh
distance, the surface no longer behaves as a point re-radiator but as a mirror (or lens), and the
path loss approaches that of specular reflection, $\propto1/(d_1+d_2)^2$; measurements by Tang and
colleagues (2021) confirmed both regimes. In that regime the RIS is as good as a perfect mirror,
which is the most it can ever be.

\mdcfig{ch25_ris}{(a) Received power through an optimally phased RIS 95\,m from a base station and
7\,m from the user ($P_t=20$\,dBm, 15\,dBi base-station antenna, free space), against the number of
half-wavelength elements, compared with the line-of-sight direct path. The RIS path grows 20\,dB per
decade of $N$. (b) Elements needed for the RIS path to equal the direct path, from
\eqref{eq:ch25:riseq}, as the RIS slides along a 100\,m link at 5\,m from the line: hundreds near the
ends at 3\,GHz, thousands at 28\,GHz, tens of thousands at 140\,GHz.}

\begin{worked}[How large must an RIS be to beat the direct path?]
A 28\,GHz base station serves a user 100\,m away; an RIS can be mounted on a wall 95\,m from the base
station and 7.07\,m from the user. At 28\,GHz, $\lambda=1.07$\,cm and a half-wavelength element has
$A=(\lambda/2)^2=2.87\times10^{-5}$\,m$^2$. From~\eqref{eq:ch25:riseq}, the RIS path equals an unobstructed
direct path when
\[
NA=\frac{0.0107\times95.1\times7.07}{100}=0.072\ \text{m}^2,\qquad N=\frac{0.072}{2.87\times10^{-5}}\approx2500,
\]
a panel about 27\,cm square. With the direct path present, the coherent sum then gives at most
$|1+1|^2$, a 6\,dB gain. The RIS earns its keep when the direct path is \emph{blocked}: with 20\,dB of
blockage, a panel with one tenth of that area, about 250 elements (8.5\,cm square), restores the
link to the blocked level and roughly doubles the received amplitude; to restore the unblocked level
takes the full 2500. With 2-bit phase control the coherent gain loses about 0.9\,dB, with 1-bit about
3.9\,dB (the loss is $20\log_{10}\sinc(2^{-b})$ for $b$-bit uniform quantisation). Repeat the
calculation at 3.5\,GHz and the same 27\,cm panel holds only 40 elements, yet the required number of
elements falls by $28/3.5=8$ to about 310: at lower frequencies an RIS must be physically much larger.
\end{worked}

\subsubsection{RIS versus relays and repeaters}
A passive RIS adds no noise and works in full duplex, but it cannot amplify. A decode-and-forward
relay amplifies and regenerates, but a half-duplex relay spends half the time listening. Bj\"ornson,
\"Ozdogan and Larsson (2020) compared them at the same location and found that an RIS needs hundreds
to thousands of elements to beat a relay at the low and moderate rates typical of the cell edge, and
wins at high rates where the half-duplex penalty dominates. In practice, the competitor that reached
standards first was the \emph{network-controlled repeater} of 3GPP Release~18: an amplify-and-forward
repeater with beamforming at both sides, configured by the base station. It needs power but no complex
channel estimation.

\begin{inpractice}[What stands between RIS and deployment]
An RIS must be told how to configure itself, so it needs a control link and someone to compute the
configuration from channel knowledge it cannot measure itself: estimating $N$ cascaded channel
coefficients with passive elements is a pilot-overhead problem that grows with $N$. Elements based on
PIN diodes or varactors draw small but non-zero power, and a large panel with fast switching needs a
supply. An RIS also reflects \emph{everything} in its bandwidth, including other operators' signals,
which raises regulatory questions about who controls a surface that redirects someone else's
spectrum. ETSI's Industry Specification Group on RIS (created in 2021) has published group reports on
use cases, channel models and architecture. The near-term applications are niche but real: coverage
of fixed blind spots in mmWave networks, transparent window films that pass or redirect particular
bands, and indoor factory deployments.
\end{inpractice}

\subsection{Fluid and movable antennas}

A different way to add spatial degrees of freedom is to move the antenna. A \emph{fluid antenna}
(Wong and colleagues, 2021) is a conductive liquid or a reconfigurable pixel structure whose radiating
position can be switched among many ports within a small space; a \emph{movable antenna} (Zhu, Ma and
Zhang, 2023) is mechanically repositioned by a small actuator. Either can pick, within a few
wavelengths, the position where the fading is favourable---selection diversity in space with one RF
chain---or where the interference is in a null. The idea is old (it is selection diversity, Chapter~\ref{ch:mimo},
with many closely spaced candidate positions), the analysis is new, and the practical questions are the speed of switching,
the accuracy of the channel knowledge needed to choose, and mechanical reliability. Treat it as a
promising research topic rather than a 6G feature.

%======================================================================
\section{Integrated sensing and communication}
\label{sec:ch25:isac}
%======================================================================

Radar and communication grew from the same root---Hertz's spark gaps, then the vacuum tube, then the
cavity magnetron of 1940---and then grew apart into separate industries, bands and standards. They
are converging again. A modern base station has wide bandwidth, a large array, accurate
synchronisation and plenty of processing, which is exactly the description of a radar.
\term{Integrated sensing and communication} (ISAC, also called joint communication and sensing) makes
the network sense objects that carry no device: cars and drones, pedestrians, rain, the layout of a
room, a breathing patient. IMT-2030 lists it as one of its six usage scenarios, 3GPP SA1 studied its
use cases in Release~19 (TR~22.837), and RAN extended its channel model (TR~38.901) to include
targets and background clutter for sensing evaluation.

\subsection{Sensing modes and use cases}

Three geometries are possible. In \emph{monostatic} sensing a base station or device transmits and
receives its own echoes, like a classic radar; it knows its own waveform exactly, but it must receive
a weak echo while transmitting, which requires full-duplex isolation of 100\,dB or more between the
transmit and receive chains, or gaps in transmission. In \emph{bistatic} or \emph{multistatic}
sensing, one node transmits and others receive; there is no self-interference problem, but the
receivers must know the transmitter's timing, frequency and position precisely, which the cellular
network's synchronisation can provide. In \emph{device-based} sensing, a phone or vehicle uses the
downlink signals as illumination. Use cases studied for 6G include intrusion and drone detection
around critical sites, vehicle and pedestrian tracking at intersections, rail-track monitoring,
indoor presence and gesture detection, rainfall estimation from link attenuation, and building a
continuously updated map of the radio environment that improves beam management.

\subsection{OFDM radar processing}

Chapter~\ref{ch:ofdm} showed that an OFDM frame is a good radar waveform. Recall the result. A target
at range $R$ with radial velocity $v$ returns the transmitted grid $X[m,n]$ (subcarrier $m$, symbol
$n$) as
\begin{equation}
Y[m,n]=a\,X[m,n]\,e^{-j2\pi m\Delta f\cdot2R/c}\,e^{\,j2\pi nT_s\cdot2v/\lambda}+W[m,n],
\label{eq:ch25:ofdmradar}
\end{equation}
where $T_s$ is the symbol duration including the cyclic prefix. Dividing by the known $X$ removes the
data and leaves a two-dimensional complex sinusoid; an inverse DFT over subcarriers gives range, and a
DFT over symbols gives velocity. The resolutions and unambiguous limits are
\begin{equation}
\Delta R=\frac{c}{2M\Delta f}=\frac{c}{2B},\quad
R_{\max}=\frac{c}{2\Delta f},\quad
\Delta v=\frac{\lambda}{2NT_s},\quad
v_{\max}=\pm\frac{\lambda}{4T_s}.
\label{eq:ch25:radarres}
\end{equation}
The 2-D DFT also provides a processing gain of $MN$, so a target can be detected even when its echo
is below the noise in every resource element. Resolution is not accuracy: with enough SNR, a peak can
be located to a small fraction of a bin. The Cram\'er--Rao bound for the delay of a known signal with
RMS bandwidth $\beta_{\text{rms}}$ and post-processing SNR $\rho$ is
$\operatorname{var}(\hat\tau)\ge1/(8\pi^2\beta_{\text{rms}}^2\rho)$, so range accuracy improves with
both bandwidth and SNR, while resolution---the ability to separate two targets---depends on bandwidth
alone.

\begin{worked}[An ISAC range--Doppler budget at 28\,GHz with 400\,MHz]
Use NR numerology $\mu=3$: $\Delta f=120$\,kHz, $T_s\approx8.93\,\mu$s including the normal cyclic
prefix of about 0.57\,$\mu$s. Take 3300 subcarriers, the 275 resource blocks that are the most an NR
carrier can hold (a standard 400\,MHz channel actually configures 264 blocks, 3168 subcarriers, which
changes the numbers below by only 4\%), so $B=3300\times120\,\text{kHz}=396$\,MHz.
\begin{itemize}
\item Range resolution: $\Delta R=3\times10^8/(2\times396\times10^6)=0.38$\,m. (With exactly
400\,MHz, 0.375\,m.) Two cars a metre apart are resolved; a window that lowers the sidelobes widens
the main lobe by roughly a factor of two, to about 0.75\,m.
\item Unambiguous range: $c/(2\Delta f)=1250$\,m. But a target is ISI-free only if its round-trip delay
fits in the cyclic prefix: $cT_{\text{cp}}/2=3\times10^8\times0.57\times10^{-6}/2\approx85$\,m. Longer
ranges need processing that tolerates the inter-symbol leakage, or a longer prefix.
\item Velocity with $N=256$ symbols (2.3\,ms): $\lambda=1.07$\,cm,
$\Delta v=0.0107/(2\times2.29\times10^{-3})=2.3$\,m/s, and $v_{\max}=\pm\lambda/(4T_s)=\pm300$\,m/s.
\item Processing gain: $10\log_{10}(3300\times256)=59$\,dB.
\end{itemize}
Figure~\ref{fig:ch25_isac_rd} shows the resulting map with five targets, including two cars at the
same velocity one metre apart, which the 396\,MHz waveform separates and a 99\,MHz waveform does not.
\end{worked}

\mdcfig{ch25_isac_rd}{OFDM radar with the parameters of the worked example (28\,GHz, 3300 subcarriers
at 120\,kHz, 256 symbols, QPSK data, Hann windows). (a) Range--Doppler map of five targets with a per
resource-element SNR of $-15$\,dB for the strongest: a pedestrian at 12\,m, a parked car at 48\,m,
two cars at 35 and 36\,m moving at 15\,m/s, and an oncoming car at 60\,m. (b) Range profiles through
the 15\,m/s column: with 396\,MHz the two cars are resolved; with 99\,MHz they merge.}

\subsection{The sensing--communication trade-off}

If an OFDM frame is already a good radar waveform, what is the trade-off? There are several. Time,
frequency, power and beams spent on sensing are not spent on data: a beam that illuminates an
intersection to track cars does not serve a user elsewhere. Monostatic sensing needs listening
time or full duplex. And, more subtly, the \emph{statistics} of the transmitted symbols matter.

Suppose the radar uses a correlation (matched-filter) receiver, multiplying $Y[m,n]$ by $X^*[m,n]$
instead of dividing by $X[m,n]$. For a target in range bin 0, the range profile after the IDFT is
\begin{equation}
z[\ell]=\frac{a}{M}\sum_{m=0}^{M-1}|X[m]|^2\,e^{\,j2\pi m\ell/M}.
\end{equation}
If the symbols have constant modulus, $|X[m]|^2=1$ and the sum is $M$ at $\ell=0$ and zero elsewhere:
perfect sidelobes. If the amplitudes fluctuate, the fluctuation $|X[m]|^2-1$ is a random sequence whose
transform spreads over all range bins. Its power relative to the peak is
\begin{equation}
\frac{\E\big[|z[\ell\ne0]|^2\big]}{|z[0]|^2}\approx\frac{\kappa-1}{M},\qquad\kappa=\frac{\E|X|^4}{(\E|X|^2)^2},
\label{eq:ch25:floor}
\end{equation}
where $\kappa$ is the \emph{kurtosis} of the constellation: 1 for PSK, 1.32 for 16-QAM, 1.38 for 64-QAM
and 2 for complex Gaussian symbols. A weak target in a sidelobe floor 30\,dB below a strong one is
lost (Figure~\ref{fig:ch25_isac_tradeoff}a). Dividing by $X$ (the reciprocal filter) removes the floor
but amplifies the noise on subcarriers that carried low-amplitude symbols by $\E[1/|X|^2]$, which is
finite for QAM and infinite for Gaussian signalling.

Communication, on the other hand, wants exactly the opposite: capacity on the Gaussian channel is
achieved by Gaussian, high-kurtosis inputs, and at 15\,dB a constant-modulus 16-PSK carries about
0.4 bits per symbol less than 16-QAM. Xiong and colleagues (2023) formalised this as the
\emph{deterministic--random trade-off}: sensing performance is best with deterministic,
constant-envelope signals, communication with random, Gaussian-like ones, and the Pareto boundary
between the two is traced by input distributions that interpolate between them
(Figure~\ref{fig:ch25_isac_tradeoff}b).

\mdcfig{ch25_isac_tradeoff}{The deterministic--random trade-off. (a) Range profile of a single target
with a correlation receiver and $M=1024$ subcarriers, averaged over 30 frames: constant-modulus 16-PSK
has no data-induced floor; 64-QAM and Gaussian symbols raise it to $(\kappa-1)/M$. (b) Each
constellation as a point: communication performance (mutual information at 15\,dB SNR) against the
sensing penalty (data-induced sidelobe floor). Better communication costs sensing quality.}

\begin{keyidea}[Sensing is estimation; communication is detection]
A communication receiver must decide among a finite number of messages and does not care what the
channel was, as long as it is equalised. A sensing receiver wants exactly that channel---the delays,
Dopplers and angles---and treats the data as a nuisance. The same signal serves both, but the
optimal design for each is different, and every ISAC system is a compromise along the trade-off of
Figure~\ref{fig:ch25_isac_tradeoff}. In practice the compromise is usually made with resources (some
symbols, beams or reference signals dedicated to sensing) rather than with the constellation.
\end{keyidea}

\begin{pitfall}[A network that sees everyone]
A cellular network that can track devices it does not serve, through walls, is also a surveillance
system. Sensing data about people is personal data under most privacy laws, and standards bodies have
made privacy and the authorisation of sensing requests part of the requirements. An engineer designing
ISAC should treat ``what may this network sense, for whom, and with whose consent'' as a design
question, not an afterthought.
\end{pitfall}

%======================================================================
\section{An AI-native air interface}
\label{sec:ch25:ai}
%======================================================================

Machine learning entered wireless networks long before ``AI-native'' became a slogan: in traffic
prediction, self-organising networks, fault detection and energy management. What is new is its
entry into the \emph{air interface}, the part defined by the physical-layer standard, where a learned
component in one device must interoperate with components built by other vendors.

\subsection{Three levels of ambition}

It helps to separate three levels.
\begin{enumerate}
\item \textbf{Learning inside one box.} A vendor replaces a block of its own receiver---a channel
estimator, a detector, a decoder schedule---with a trained model. Nothing changes over the air, and no
standardisation is needed beyond performance requirements. This is already common.
\item \textbf{Learning across the interface.} The model on one side produces something the other side
must understand: compressed channel-state feedback, predicted beams, positioning measurements. The
standard must define the interface, the data used for training, how models are identified, activated,
monitored and updated (the model \emph{life cycle}), and how performance is tested. 3GPP's Release~18
study (TR~38.843) examined three use cases---CSI feedback enhancement, beam management and
positioning---and Release~19 began normative work on AI/ML-based beam management and positioning,
with CSI compression continuing as a study because of its interoperability difficulty.
\item \textbf{Learning the air interface itself.} Transmitter and receiver are trained jointly,
end to end, so that the waveform, constellation and pilots are learned rather than designed. This is
the research frontier.
\end{enumerate}

\subsection{End-to-end learning: the autoencoder view}

O'Shea and Hoydis (2017) pointed out that a communication system is an \emph{autoencoder}: a
transmitter $f_{\bm{\theta}}$ maps a message $s\in\{1,\dots,M\}$ to a channel input
$\vect{x}=f_{\bm{\theta}}(s)$, the channel adds a ``bottleneck'' of noise, and a receiver
$g_{\bm{\phi}}$ maps the output to a probability vector over the messages. Train both with the
cross-entropy loss
\begin{equation}
\mathcal{L}(\bm{\theta},\bm{\phi})=-\E\big[\log\,[g_{\bm{\phi}}(\vect{y})]_s\big],\qquad
\vect{y}=f_{\bm{\theta}}(s)+\vect{w},
\label{eq:ch25:autoenc}
\end{equation}
under a power constraint, by stochastic gradient descent. The gradient with respect to the
transmitter parameters flows \emph{through} the channel, which is possible when the channel model is
differentiable, as additive noise is. The cross-entropy is closely related to the mutual information:
for a receiver that outputs the true posterior it equals $H(S|\vect{Y})$, so minimising it maximises
$I(S;\vect{Y})$ for uniform messages.

Figure~\ref{fig:ch25_autoencoder} shows the result for $M=16$ messages in two real dimensions,
trained in NumPy in about a minute. The learned constellation is not square: it approximates the
densest packing of circles in a disc, with points on a roughly hexagonal lattice, and it is a fraction
of a decibel better than 16-QAM in symbol error rate. That is the \emph{packing gain} of
Chapter~\ref{ch:infotheory}, rediscovered by gradient descent. It is also a sobering result: the
network found what Foschini, Gitlin and Weinstein had found by optimisation in 1974, and the gain is
small because good constellations were already close to optimal on this channel.

\mdcfig{ch25_autoencoder}{An autoencoder trained end to end for 16 messages on the AWGN channel at
$E_s/N_0=14$\,dB (learned points with a 2-64-64-16 neural decoder; script \texttt{ch25\_figs.py}). (a)
The learned constellation and the decoder's decision regions, with 16-QAM for comparison. (b) Symbol
error rate of the learned points and of 16-QAM, both with minimum-distance detection: the learned
geometry gives a small packing gain.}

Where end-to-end learning becomes interesting is when the channel is \emph{not} well modelled:
nonlinear amplifiers, low-resolution converters, unknown interference, optical channels with
nonlinearity. There the learned transmitter can discover shaping, pilot patterns or pulse shapes that
no textbook model suggests. The practical obstacle is that the real channel is not differentiable;
solutions include learning a differentiable surrogate (a generative model of the channel), training
the transmitter with reinforcement-learning-style gradient estimates (Aoudia and Hoydis, 2019), or
fixing the transmitter and learning only the receiver.

\subsection{Neural receivers}

A \emph{neural receiver} replaces the chain of channel estimation, equalisation and demapping in an
OFDM receiver with a single convolutional network that takes the received resource grid (data and
pilots) and outputs bit log-likelihood ratios for the channel decoder. Published examples such as
DeepRx (Honkala, Korpi and Huttunen, 2021) and the receivers of Aoudia and Hoydis (2022) outperform
conventional least-squares estimation with linear interpolation, particularly with sparse pilots and
high Doppler, and can make pilots partly unnecessary by exploiting a learned, slightly non-uniform
constellation (``pilotless'' transmission). Open tools such as NVIDIA's Sionna library make such
experiments routine. The caveats are complexity (a convolutional network over a grid of thousands of
resource elements costs far more operations than a linear MMSE estimator), generalisation to
conditions not seen in training, and the difficulty of guaranteeing worst-case performance.

\subsection{CSI compression}

Type~II CSI feedback in NR (Chapter~\ref{ch:mimo}) describes a channel as a weighted sum of a few DFT
beams. A neural autoencoder can do better: an encoder in the terminal compresses the measured
channel to a short code word; a decoder in the base station reconstructs it. CsiNet (Wen, Shih and Jin,
2018) introduced the idea; many variants followed. The gain comes from learning the \emph{statistics}
of real channels in a particular deployment, which a fixed codebook ignores. The difficulty is that
the encoder and decoder are built by different companies. If the terminal's encoder was trained with
one vendor's decoder, will it work with another's? Options studied in 3GPP include jointly trained
reference models, sequential training (one side trains to match the other's published outputs) and
standardised model structures, each with trade-offs between performance, proprietary freedom and
testability.

\subsection{Learned constellation shaping}

Chapter~\ref{ch:infotheory} showed that uniform QAM leaves up to 1.53\,dB of \emph{shaping gain}
unclaimed, recoverable by geometric shaping (non-uniform point positions) or probabilistic shaping
(non-uniform point probabilities, as in modern coherent optical transceivers). Autoencoders can learn
both jointly, including the bit labelling, for channels where analytic optimisation is awkward
(Stark, Ait Aoudia and Hoydis, 2019). This is perhaps the most likely route for learned PHY components
into products: a learned constellation is a fixed table once trained, cheap to implement and easy to
test.

\begin{pitfall}[The limits of learned physical layers]
A trained model is only as good as its training data. Channel conditions, interference and hardware
outside the training distribution can degrade it silently, and unlike a model-based receiver it gives
no easy way to bound its worst case. It can be attacked with crafted adversarial perturbations. Its
inference costs energy that a simple algorithm does not. And a standard must specify how two devices
from different vendors prove that they interoperate. None of these is a reason to avoid learning; all
are reasons to compare every learned component with the \emph{best} model-based baseline---not a
straw man---and to keep a fallback.
\end{pitfall}

\begin{keyidea}[Learn what you cannot model]
Fifty years of communication theory have produced near-optimal solutions for linear, Gaussian,
well-modelled channels: matched filters, MMSE estimators, LDPC codes. Machine learning rarely beats
them there. It earns its place where the model is poor or the optimal algorithm is intractable:
hardware nonlinearities, interference with structure, site-specific channel statistics, the joint
optimisation of many interacting parameters, and resource management decisions too complex for
closed-form rules.
\end{keyidea}

%======================================================================
\section{Networks in three dimensions}
\label{sec:ch25:ntn}
%======================================================================

Chapter~\ref{ch:satellite} described how 3GPP adapted 5G to satellites in Releases~17--19:
transparent and regenerative payloads, terminal-side timing advance and Doppler pre-compensation from
GNSS and ephemeris, longer HARQ timelines, and direct-to-device service to ordinary phones. 6G aims
to go further, designing the air interface from the start for a \emph{three-dimensional network} of
terrestrial cells, high-altitude platforms, low- and medium-orbit satellites and geostationary
satellites, with a terminal moving seamlessly among them (Figure~\ref{fig:ch25_3d}).

\begin{figure}[htbp]\centering
\begin{tikzpicture}[x=1cm,y=1cm,
  lay/.style={draw=navy!50,dashed},
  lab/.style={font=\sffamily\scriptsize,anchor=west,text=navy},
  nd/.style={draw=navy,fill=navy!8,thick,rounded corners=2pt,font=\sffamily\tiny,align=center,inner sep=2pt},
  sat/.style={draw=accent,fill=accent!10,thick,rounded corners=2pt,font=\sffamily\tiny,align=center,inner sep=2pt},
  lk/.style={thick,navy!70},
  isl/.style={thick,accent,dashed}]
\fill[green!15!white] (-0.2,0) rectangle (11.2,0.35);
\draw[navy!60] (-0.2,0.35)--(11.2,0.35);
\node[lab] at (11.25,0.2) {ground};
\draw[lay] (-0.2,1.5)--(11.2,1.5); \node[lab] at (11.25,1.5) {HAPS, $\sim$20\,km};
\draw[lay] (-0.2,3.0)--(11.2,3.0); \node[lab] at (11.25,3.0) {LEO, 500--1500\,km};
\draw[lay] (-0.2,4.2)--(11.2,4.2); \node[lab] at (11.25,4.2) {GEO, 35\,786\,km};
\node[nd] (bs1) at (1.0,0.75) {macro\\gNB};
\node[nd] (bs2) at (3.3,0.75) {small\\cell};
\node[nd] (ue1) at (2.2,0.6) {UE};
\node[nd] (ue2) at (5.6,0.6) {UE};
\node[nd] (ue3) at (8.3,0.6) {UE};
\node[nd] (ue4) at (10.0,0.6) {IoT};
\node[nd] (gw) at (7.0,0.75) {gateway};
\node[sat] (haps) at (5.0,1.5) {HAPS\\(HIBS)};
\node[sat] (leo1) at (6.3,3.0) {LEO};
\node[sat] (leo2) at (8.8,3.0) {LEO};
\node[sat] (geo) at (9.8,4.2) {GEO};
\node[nd] (core) at (1.0,2.4) {6G core\\+ edge};
\draw[lk] (bs1)--(ue1); \draw[lk] (bs2)--(ue1);
\draw[lk] (haps)--(ue2); \draw[lk] (haps)--(bs2);
\draw[lk] (leo1)--(ue3); \draw[lk] (leo1)--(gw);
\draw[lk] (leo2)--(ue4); \draw[isl] (leo1)--node[above,font=\sffamily\tiny,text=accent]{inter-satellite link}(leo2);
\draw[lk] (geo)--(leo2);
\draw[lk,dotted] (core)--(bs1); \draw[lk,dotted] (core)--(gw);
\end{tikzpicture}
\caption{A three-dimensional network. Terrestrial macro and small cells carry most traffic;
high-altitude platforms (HAPS) acting as IMT base stations (HIBS) fill rural and disaster areas;
LEO satellites with inter-satellite links serve terminals directly and backhaul remote cells; GEO
satellites provide broadcast and wide-area coverage. The 6G design goal is one air interface and one
core network for all of them.}
\label{fig:ch25_3d}
\end{figure}

\paragraph{High-altitude platforms.} A \term{high-altitude platform station} (HAPS) is an aircraft,
airship or balloon in the stratosphere at about 20\,km, above the weather and commercial air
traffic. At 20\,km a platform sees the ground down to 15$^\circ$ elevation out to a radius of roughly
$20/\tan15^\circ\approx75$\,km, a footprint of over 17\,000\,km$^2$, with a round-trip delay below
0.5\,ms and a free-space loss only a few decibels worse than a large terrestrial macro cell at its
edge. WRC-23 identified certain IMT bands below 2.7\,GHz for HAPS acting as IMT base stations
(HIBS), so a HAPS can serve unmodified phones. The engineering challenges are endurance (solar-powered
aircraft must survive the night, especially in winter at high latitudes), station keeping in
stratospheric winds, and payload power and weight. Several programmes have flown demonstrators;
none has yet operated a commercial service at scale.

\paragraph{One system, many altitudes.} The architectural ideas for 6G non-terrestrial integration
include a common air interface for terrestrial and satellite access with satellite-specific
parameters (longer prefixes, larger HARQ process counts, Doppler pre-compensation) as configurable
options; regenerative payloads with on-board base-station functions and inter-satellite links, so
that traffic can be routed in space; mobility management that predicts satellite passes from
ephemeris; and spectrum sharing between terrestrial and satellite layers in the same or adjacent
bands. Direct-to-device service, which in 2024--2026 meant text messaging and low-rate data to phones,
is expected to grow towards broadband as satellite apertures grow (the link budget in
Chapter~\ref{ch:satellite} showed why the uplink sets the aperture).

%======================================================================
\section{Semantic and goal-oriented communication}
\label{sec:ch25:semantic}
%======================================================================

In 1949 Warren Weaver, introducing Shannon's theory to a general audience, distinguished three levels
of the communication problem: the \emph{technical} problem (how accurately can the symbols be
transmitted?), the \emph{semantic} problem (how precisely do the transmitted symbols convey the desired
meaning?) and the \emph{effectiveness} problem (how effectively does the received meaning affect conduct
in the desired way?). Shannon solved level A, explicitly setting meaning aside
(Chapter~\ref{ch:infotheory}). \term{Semantic communication} proposes to engineer levels B and C:
transmit only what is needed for the receiver to reconstruct the meaning, or to accomplish a task.

The most concrete form is \emph{deep joint source--channel coding} (deep JSCC): a neural encoder maps
an image directly to channel symbols, and a decoder reconstructs it, trained end to end
(Bourtsoulatze, Kurka and G\"und\"uz, 2019). Unlike a separated design (a JPEG encoder followed by an
LDPC code), its quality degrades gracefully as the SNR falls instead of collapsing at the code's
threshold---the ``cliff effect''---which matters for broadcasting to receivers with different SNRs and
for low latency. \emph{Goal-oriented} communication goes further: a camera on a robot need not send
the image, only what the controller needs, perhaps a few bytes of object positions, and the right
metric is task performance or the \emph{age of information} (how fresh the receiver's knowledge is),
not bit error rate.

A critical view is essential. The source--channel separation theorem
(Chapter~\ref{ch:infotheory}) says that, for long blocks, point-to-point links and a fixed distortion
measure, separate source and channel coding is optimal. Semantic gains therefore come from where the
theorem's assumptions fail: short blocks and tight latency, broadcast to heterogeneous receivers,
and---most of all---a better \emph{source model}, which is just good compression. Sending object
positions instead of video is application design; it does not need a new physical layer. And
``meaning'' is not yet a quantity with a theory comparable to Shannon's; many papers measure semantic
performance with task metrics specific to their own datasets. The likely legacy of semantic
communication is a set of useful techniques---learned JSCC for some media, task-aware compression,
freshness-aware scheduling---rather than a replacement of the bit as the unit of communication.

\begin{keyidea}[Shannon's separation still sets the bar]
Whenever a proposal claims to beat ``conventional'' communication by exploiting meaning, ask: is the
gain from a better source model (which a separated system could also use), from finite-blocklength or
broadcast effects (real but bounded), or from changing the task (an application decision)? The
answers are often all three, and the size of each tells you how much of the idea belongs in the air
interface.
\end{keyidea}

%======================================================================
\section{Energy efficiency and sustainability}
\label{sec:ch25:energy}
%======================================================================

Mobile networks consume a substantial share of an operator's operating costs in electricity, and the
radio access network---the base stations---accounts for most of a mobile network's energy use.
Two metrics are common. \term{Energy efficiency} in bits per joule (or its inverse, joules per bit)
measures how much data is delivered per unit energy; it improved by orders of magnitude from 2G to
5G as spectral efficiency and bandwidth grew faster than power consumption. \emph{Absolute}
consumption is what appears on the electricity bill and in emissions inventories, and it has not
fallen, because traffic has grown even faster.

The difficulty is that base stations are not \emph{energy-proportional}: they consume a large fraction
of their peak power when idle. A widely used model, from the European EARTH project (Auer and
colleagues, 2011), writes the input power of a site with $N_{\text{TRX}}$ transceiver chains as
\begin{equation}
P_{\text{in}}=N_{\text{TRX}}\left(P_0+\Delta_pP_{\text{out}}\right)\quad(0<P_{\text{out}}\le P_{\max}),
\qquad P_{\text{in}}=N_{\text{TRX}}P_{\text{sleep}}\quad(P_{\text{out}}=0),
\label{eq:ch25:earth}
\end{equation}
with, for a macro site of the LTE era, $P_0\approx130$\,W, $\Delta_p\approx4.7$, $P_{\max}=20$\,W and
$P_{\text{sleep}}\approx75$\,W per chain. At zero load the site still draws 58\% of its full-load power
(Figure~\ref{fig:ch25_energy}a). Traffic follows a daily cycle with a deep trough at night, so much
of the energy is spent waiting.

\mdcfig{ch25_energy}{(a) Input power of a three-sector macro site with six transceiver chains against
load, from the EARTH model~\eqref{eq:ch25:earth}: always on, with sleep in idle symbols, and the ideal
of power proportional to load. (b) One day with an illustrative traffic profile: sleeping in empty
symbols saves about a fifth of the daily energy here; deeper sleep modes and switching off carriers
at night save more.}

The levers are therefore at least as much about \emph{doing nothing efficiently} as about doing
something efficiently:
\begin{itemize}
\item \textbf{Lean carriers.} LTE transmitted cell-specific reference signals in every subframe,
whether or not anyone was listening. NR removed them and made synchronisation blocks sparse
(every 20\,ms by default), so a base station can sleep between them. 6G aims to go further, with
on-demand system information and synchronisation signals that can be switched off in cells with no
users.
\item \textbf{Sleep modes.} Components can be switched off in idle symbols (micro-sleep, below a
millisecond), idle subframes, or longer periods (deep sleep, carrier and cell switch-off), each with
a longer wake-up time and a larger saving. 3GPP's Release~18 study on network energy savings (TR~38.864)
and the resulting features added mechanisms for cell discontinuous transmission and reception and for
adapting the number of active antennas and the bandwidth to the load.
\item \textbf{Hardware.} Power amplifier efficiency (Doherty and envelope-tracking designs,
digital predistortion), GaN devices, and silicon process scaling for the baseband.
\item \textbf{Network design.} Fewer, better-placed sites; offloading traffic to the most efficient
layer at night; and cooling, which in some sites consumes as much as the radio.
\end{itemize}

\begin{worked}[Joules per bit at a macro site]
A site draws 1.0\,kW at an average load that delivers 300\,Mb/s across its sectors. Its energy per bit
is $1000/(300\times10^6)=3.3\,\mu$J/bit. At night, the load falls by a factor of ten but, from
Figure~\ref{fig:ch25_energy}, the power falls only to about 0.8\,kW: $0.8\times10^3/(30\times10^6)=27\,
\mu$J/bit, eight times worse. A phone downloading at 100\,Mb/s with a modem drawing 1\,W spends
10\,nJ/bit---the network side dominates by orders of magnitude, which is why network energy saving,
not handset efficiency, is the focus of the sustainability work in 6G.
\end{worked}

Sustainability is broader than electricity. The embodied energy and materials of equipment, the
lifetime of devices, and the ``enabling'' effect of connectivity on other sectors (remote work,
logistics, smart grids) are all part of the assessment, and all are hard to measure. IMT-2030 lists
sustainability as an overarching aspect; turning it into testable requirements is part of the work
in progress.

%======================================================================
\section{Zero-energy devices and backscatter}
\label{sec:ch25:ambient}
%======================================================================

At the opposite end of the power scale are devices with no battery at all. Chapter~\ref{ch:wifi}
described passive RFID: a reader's continuous wave powers the tag, which answers by switching its
antenna load and so modulating the reflected wave---\term{backscatter}. Ambient backscatter
(Liu and colleagues, 2013) showed that a tag can instead modulate signals that are already in the air,
such as TV or cellular transmissions, and that a nearby receiver can decode the tiny modulation by
comparing the signal with and without the tag's reflection.

3GPP studied \term{ambient IoT} in Release~19 (TR~38.848) and began specifying it: devices that harvest
energy from radio waves, light, vibration or temperature differences, consuming microwatts, reached
by a base station, a UE or a dedicated reader. The target applications are logistics (tracking
pallets and parcels through warehouses), asset inventory and sensors embedded in structures, where
the cost of a battery or of replacing it exceeds the value of the data.

The link budget of a backscatter device is the RIS product-distance law in miniature. The power
reaching the device falls as $1/d_1^2$, and the reflected power at a receiver at distance $d_2$
falls as $1/(d_1^2d_2^2)$. Two limits follow. The \emph{forward} link must deliver enough power for
the device's energy harvester, typically somewhere around $-20$ to $-30$\,dBm for simple tags, which
limits the reader-to-device distance to metres or tens of metres. The \emph{backward} link, being a
double path loss, is weak, but receivers can integrate over time because the data rates are low. The
6G question is whether a cellular network can make such devices reachable everywhere, without
dedicated readers.

%======================================================================
\section{Security in the 6G era}
\label{sec:ch25:security}
%======================================================================

\subsection{Physical-layer security}

Chapter~\ref{ch:infotheory} introduced the wiretap channel: if the legitimate receiver's channel is
better than the eavesdropper's, codes exist that deliver a positive secrecy rate,
$C_s=[\log_2(1+\mathrm{SNR}_B)-\log_2(1+\mathrm{SNR}_E)]^+$, with no shared key. Large arrays, beamfocusing
and RIS can, in principle, make Bob's channel much better than Eve's. A related and more practical
idea is \emph{secret-key generation}: because the channel between two devices is reciprocal and
changes randomly, both ends can quantise their channel measurements into matching bit strings that an
eavesdropper more than half a wavelength away cannot reproduce. Both ideas face the same obstacle:
the system rarely knows where Eve is or how good her receiver is, so the guarantee depends on
assumptions an attacker need not respect. Physical-layer security is best viewed as an additional
layer---for example for lightweight IoT devices or for key refresh---rather than a replacement for
cryptography.

\subsection{The post-quantum migration}

A more certain change is the migration to \term{post-quantum cryptography}. A large fault-tolerant
quantum computer running Shor's algorithm would break the public-key systems in use today---RSA and
elliptic-curve cryptography---in polynomial time. Nobody knows when, or whether, such a machine will be
built, but the threat is already relevant because of ``harvest now, decrypt later'': traffic recorded
today could be decrypted years from now. In August 2024 NIST published the first post-quantum standards:
FIPS~203 (ML-KEM, a lattice-based key-encapsulation mechanism derived from CRYSTALS-Kyber),
FIPS~204 (ML-DSA, lattice-based signatures) and FIPS~205 (SLH-DSA, hash-based signatures).

For mobile networks the exposure is concentrated in a few places. 5G encrypts the subscriber's
permanent identity with an elliptic-curve scheme (ECIES with Curve25519 or secp256r1, defined in TS~33.501)
before sending it over the air; network functions authenticate each other with TLS and certificates;
and operators connect over IPsec. All of these must move to quantum-resistant or hybrid schemes. The
symmetric algorithms that protect user data are less affected: Grover's algorithm only halves the
effective key length, so moving from 128-bit to 256-bit keys suffices, and 3GPP studied 256-bit
algorithms for 5G (TR~33.841). Post-quantum algorithms have a communication-engineering cost: an ML-KEM-768
public key is 1184 bytes and its ciphertext 1088 bytes, against 32 bytes for an X25519 key, which
matters for messages that must fit in a single radio transport block or a constrained IoT frame.
6G is expected to be designed with quantum-resistant algorithms from the start and, just as
important, with \emph{crypto-agility}---the ability to replace algorithms without replacing the
system.

%======================================================================
\section{Quantum communication, briefly}
\label{sec:ch25:quantum}
%======================================================================

Quantum mechanics offers a different route to secure keys. In \term{quantum key distribution} (QKD),
the security rests not on the difficulty of a mathematical problem but on physics: measuring a quantum
state disturbs it, and an unknown quantum state cannot be copied (the no-cloning theorem).

\subsection{BB84 in five steps}
Bennett and Brassard's 1984 protocol, BB84, uses single photons polarised in one of two
\emph{bases}: rectilinear ($+$: horizontal for 0, vertical for 1) or diagonal ($\times$: 45$^\circ$ for
0, 135$^\circ$ for 1).
\begin{enumerate}
\item Alice sends a stream of photons, choosing for each a random bit and a random basis.
\item Bob measures each photon in a randomly chosen basis. If his basis matches Alice's he gets her bit;
if not, his result is random.
\item Over a public, authenticated classical channel, they compare \emph{bases} (not bits) and keep
only the positions where they agree, about half: the \emph{sifted key}.
\item They sacrifice a random sample of the sifted key to estimate the quantum bit error rate (QBER).
An eavesdropper who measures and resends photons must guess bases and introduces errors at a rate of
25\% in the bits she intercepts. If the QBER is below a threshold (about 11\% for BB84 with one-way
processing, from the security proof of Shor and Preskill, 2000), they proceed.
\item They correct the remaining errors with a classical error-correcting code
(\emph{information reconciliation}, often with LDPC codes, Chapter~\ref{ch:moderncodes}) and shorten
the key with a hash function to remove any information Eve could have (\emph{privacy amplification}).
\end{enumerate}

\begin{table}[htbp]\centering\small
\caption{A BB84 exchange of eight photons. Positions where the bases agree form the sifted key.}
\label{tab:ch25:bb84}
\begin{tabular}{@{}l *{8}{c}@{}}
\toprule
Alice's bit & 0 & 1 & 1 & 0 & 1 & 0 & 0 & 1\\
Alice's basis & $+$ & $\times$ & $+$ & $+$ & $\times$ & $\times$ & $+$ & $\times$\\
Bob's basis & $+$ & $+$ & $+$ & $\times$ & $\times$ & $+$ & $+$ & $\times$\\
Bob's result & 0 & (rand.) & 1 & (rand.) & 1 & (rand.) & 0 & 1\\
Kept? & \checkmark & & \checkmark & & \checkmark & & \checkmark & \checkmark\\
\midrule
Sifted key & 0 & & 1 & & 1 & & 0 & 1\\
\bottomrule
\end{tabular}
\end{table}

\subsection{Distance, rate and reality}
Single photons cannot be amplified without destroying the security (an amplifier is a cloner), so loss
is the enemy. In fibre at 0.2\,dB/km, 100\,km transmits $\eta=1\%$ of the photons and 500\,km
$10^{-10}$. Pirandola and colleagues (2017) showed that no repeaterless protocol can exceed
$-\log_2(1-\eta)\approx1.44\eta$ secret bits per channel use: at a 1\,GHz pulse rate, at most about
14\,Mb/s at 100\,km but about 0.1\,b/s at 500\,km. Long distances therefore need trusted relay nodes,
protocols such as twin-field QKD whose rate scales as $\sqrt\eta$, future quantum repeaters, or
satellites: China's Micius satellite demonstrated satellite-to-ground QKD over distances up to about
1200\,km in 2017 (Liao and colleagues). QKD also needs an authenticated classical channel---which itself
needs a pre-shared key or classical signatures---and dedicated hardware, and real implementations
have been attacked through their imperfections (detector blinding, for example). National security
agencies in several countries currently recommend post-quantum cryptography rather than QKD for
general use. QKD is real and deployed in some backbone and government networks; it is not coming to
your phone.

%======================================================================
\section{Light as a radio: optical wireless}
\label{sec:ch25:owc}
%======================================================================

Visible and infrared light offer hundreds of terahertz of unlicensed spectrum, confinement by walls
(a security and reuse advantage) and no electromagnetic interference with radio systems.
\term{Optical wireless communication} (OWC) includes visible-light communication (VLC) with LEDs,
infrared links, and free-space optical (FSO) links over kilometres with lasers.

Indoor OWC almost always uses \emph{intensity modulation with direct detection} (IM/DD): the
transmitter varies the optical power, and a photodiode produces a current proportional to the received
power. The signal must be real and non-negative, so OFDM is adapted by Hermitian symmetry and a DC bias
(DCO-OFDM) or by clipping schemes (ACO-OFDM), as in Chapter~\ref{ch:ofdm}. For a line-of-sight link from
an LED with a Lambertian pattern of order $m$ to a photodiode of area $A_d$ at distance $d$, the channel
DC gain is (Kahn and Barry, 1997)
\begin{equation}
H(0)=\frac{(m+1)A_d}{2\pi d^2}\cos^m\phi\,\cos\psi,\qquad m=\frac{-\ln2}{\ln\cos\Phi_{1/2}},
\label{eq:ch25:lambert}
\end{equation}
with $\phi$ the emission angle, $\psi$ the incidence angle (within the receiver's field of view) and
$\Phi_{1/2}$ the LED's half-power semi-angle. Note the $1/d^2$ dependence of \emph{optical power}: in
IM/DD the electrical signal power at the receiver goes as the square of the optical power, so as
$1/d^4$.

The bottlenecks are the transmitters and the uplink. White phosphor-converted LEDs have modulation
bandwidths of only a few megahertz (the slow phosphor; filtering out the blue component gives tens of
megahertz); micro-LEDs and laser diodes reach hundreds of megahertz to gigahertz. Standards include IEEE
802.15.7 for short-range VLC and IEEE 802.11bb (2023), which brings light communication into the Wi-Fi
family (Chapter~\ref{ch:wifi}). In 6G, OWC is a candidate for very-high-capacity indoor hotspots,
for environments where radio is undesirable (hospitals, aircraft cabins, industrial plants), for
underwater links, and for FSO backhaul and inter-satellite links, where laser terminals are already
standard in the largest LEO constellations (Chapter~\ref{ch:satellite}).

%======================================================================
\section{What might not happen: hype and physics}
\label{sec:ch25:hype}
%======================================================================

Every generation's research phase produces more ideas than its standards can absorb. That is healthy;
it is how the good ideas are found. But an engineer must be able to tell a promising idea from a
fashionable one, and the tools for doing so are the ones in this book: the link budget, Shannon's
formula, the Rayleigh distance, the product-distance law, the Walden figure of merit, and the
economics of deployment. Table~\ref{tab:ch25:hype} applies them.

\begin{table}[htbp]\centering\small
\caption{A sceptic's guide to 6G claims. The ``assessment'' column is the author's engineering judgement
in late 2026, offered as a starting point for argument, not as a forecast.}
\label{tab:ch25:hype}
\begin{tabularx}{\linewidth}{@{}>{\raggedright\arraybackslash}p{2.7cm} >{\raggedright\arraybackslash}X >{\raggedright\arraybackslash}p{3.3cm}@{}}
\toprule
Claim & What the physics and economics say & Assessment\\
\midrule
Terahertz mobile access & Link budgets close at tens of metres with 30--40\,dBi antennas; beam tracking, body blockage and PA power are the problems & Fixed links, kiosks, data centres; not general mobile access\\
1\,Tb/s peak rates & Needs $>$100\,GHz of bandwidth or very high spectral efficiency; converter power scales with bandwidth & Laboratory and fixed links\\
RIS on every wall & Product-distance law needs thousands of elements; control and channel estimation overhead; multi-operator issues & Niche coverage fixes; repeaters first\\
OTFS/AFDM replacing OFDM & Real gains at very high Doppler; ecosystem and complexity favour OFDM & Options or precoders on OFDM\\
Fully learned air interface & Interoperability, testability and generalisation unsolved & Learned components: yes; learned PHY: not soon\\
Semantic communication & Separation theorem; gains from source models and short blocks & Techniques, not a new layer\\
0.1\,ms end to end & Light in fibre covers only 20\,km in 0.1\,ms & Air interface only; local edge\\
City-wide cell-free & Fronthaul, synchronisation and deployment cost & Clustered distributed MIMO\\
Quantum-secured phones & Loss limits, hardware, authentication & PQC everywhere; QKD in backbones\\
Network sensing & Physics sound; business model and privacy uncertain & Likely, in specific use cases\\
Upper mid-band capacity layer & Large arrays on existing sites approach mid-band coverage & Likely the main new 6G spectrum\\
\bottomrule
\end{tabularx}
\end{table}

\begin{history}[Lessons from earlier hype cycles]
The industry has been here before. The ``mobile internet'' of the WAP era around 2000 disappointed
until the smartphone arrived in 2007 and made it real, on networks designed for something else. 3G
was sold on video calling, which few people used until it became an app over 4G. WiMAX had an early
lead over LTE and lost on ecosystem. 5G millimetre-wave was promised as the defining feature of the
generation; it delivered spectacular rates where deployed but was deployed far more narrowly than early
forecasts suggested, while mid-band 5G with massive MIMO carried most of the traffic. In each case the
physics was right and the timeline, the cost or the use case was wrong. The technologies that won
were those that fitted existing sites, existing devices and an existing business.
\end{history}

\begin{pitfall}[Extrapolating the log-scale chart]
Peak data rates have grown roughly tenfold per generation, and it is tempting to draw a straight line on
a log scale to 1\,Tb/s in 6G. But the growth came from specific sources---wider channels, more antennas,
denser constellations, better codes---each of which is now near a physical or economic limit: codes are
within a fraction of a decibel of capacity, 1024-QAM needs SNRs few links achieve, and bandwidth comes
with the hardware costs of Section~\ref{sec:ch25:spectrum}. Ask which lever a forecast pulls, and how
far that lever has left to move.
\end{pitfall}

%======================================================================
\section{What endures: a closing word to the apprentice}
\label{sec:ch25:closing}
%======================================================================

We have come a long way together. We started with Chappe's semaphore towers and Morse's key, followed
the telephone and the vacuum tube to the superheterodyne radio, watched Nyquist and Shannon turn
engineering intuition into mathematics, and built up, chapter by chapter, the machinery of a modern
digital link: sampling and quantisation, filters and multirate processing, pulse shaping and matched
filtering, modulation and synchronisation, channels and equalisers, information theory and the codes
that approach its limits, OFDM, spread spectrum, MIMO, multiple access, five generations of cellular
systems, Wi-Fi, satellites and the wireline and optical networks that carry everything in the end. This
chapter has looked over the horizon. Before you close the book, let me tell you what I think will
still be true when everything in this chapter is either standard practice or forgotten.

\subsection{The ideas that do not age}

Some ideas in this book are tied to a particular generation of technology and will date. Others are
not tied to any technology at all. They are statements about signals, noise and information, and
they were as true for the telegraph as they will be for whatever replaces 6G. Carry these with you.

\begin{description}[style=nextline,leftmargin=1.2em,font=\sffamily\bfseries\small\color{navy}]
\item[Shannon's capacity, $C=B\log_2(1+\mathrm{SNR})$.] It tells you before you build anything whether a
requirement is possible, which resource---bandwidth or power---is the cheaper one to spend, and how
far a proposed scheme is from the best that can be done. Every claim in this chapter can be checked
against it. When a design is within a decibel of capacity, stop optimising the modem and look
elsewhere.
\item[The matched filter and the sufficient statistic.] Correlate with what you expect to see. It is
the optimal detector in white noise, the heart of radar and of the OFDM demodulator, the reason the
RRC filter is split, the first step of every channel estimator and of the OTFS and ISAC receivers of
this chapter. Most clever receivers are matched filters in disguise, followed by an honest treatment
of whatever the matched filter leaves behind.
\item[The link budget.] Power in, gains and losses, $kTB$ and a noise figure, and the margin that is
left. It has never once been repealed. It told you why sub-THz needs pencil beams, why direct-to-phone
satellites need enormous antennas, why an RIS must be large and why a backscatter tag cannot be far
from its reader. Write one for every system you ever touch, on paper, before you simulate.
\item[Fourier.] Time and frequency, delay and Doppler, aperture and angle, range and wavenumber: the
same transform in different clothes. OFDM, the Rayleigh distance, the radar resolutions and the
delay--Doppler domain are all Fourier pairs.
\item[Degrees of freedom.] Count them: $2BT$ dimensions in time and frequency, $\min(N_t,N_r)$ in
space, $\pi A/\lambda^2$ for an aperture. Gains come from using more of them or using them better;
anything that seems to give more than the count allows is wrong.
\item[Noise is the floor, and it is $kT$.] At 290\,K, $-174$\,dBm/Hz. You cannot cool the sky, and you
cannot amplify your way out of a noisy first stage.
\end{description}

\subsection{How to keep learning}

The half-life of specific knowledge in this field is perhaps five to ten years; the half-life of the
fundamentals is longer than a career. Spend your learning accordingly.
\begin{itemize}
\item \textbf{Read primary sources.} The original papers of Nyquist, Shannon, Viterbi, Gallager,
Berrou, Ar\i kan and Marzetta are more readable than their reputation suggests, and they show how
the ideas were found. Read the standards themselves, not summaries of them: 3GPP specifications, IEEE
802 amendments and ITU Recommendations are free, and the arguments behind them are in the
contributions.
\item \textbf{Build things.} A simulation teaches you the theory; a radio teaches you everything the
theory left out---phase noise, DC offsets, IQ imbalance, a clock that drifts, an amplifier that
clips, a neighbour's interference. The labs that accompany this book run in Python; many of them
have a ``going further'' section for a software-defined radio such as the USRP B200. Do those.
\item \textbf{Measure before you believe.} Every plot in a paper, including the ones in this book,
rests on assumptions. Reproduce the important ones yourself, change the assumptions, and see what
survives.
\item \textbf{Learn the neighbouring fields.} The breakthroughs of the last thirty years came from
crossings: coding theory and iterative decoding from statistical physics and graph theory, MIMO from
array processing and random matrices, OFDM from fast algorithms, learned receivers from machine
learning, optical coherent transmission from wireless DSP. Learn enough optimisation, statistics,
circuit design and computer architecture to talk to the people who do them.
\item \textbf{Keep a notebook of your mistakes.} Every engineer's list contains the classics: a factor
of two in a noise bandwidth, a sign error in a Doppler shift that costs a week, and a beautiful
simulation that was beautiful only because someone forgot to add the noise. You will have your own
list. Read it before every design review.
\end{itemize}

\subsection{Where the next breakthroughs might come from}

If I knew, I would be working on them. But history suggests where to look. Breakthroughs in
communication have rarely come from a new formula for capacity; they have come from making a known
limit \emph{reachable at an acceptable cost}. Turbo codes and LDPC codes did not change Shannon's
bound---they made it practical. OFDM did not change the channel---the FFT and cheap silicon made it
the obvious receiver. Massive MIMO was a theory in 2010 and a product in 2018 because digital
beamforming hardware became affordable. So look for the next one where costs are about to fall or
limits are about to bind: in transistors and packaging for the sub-THz bands, in converters that
spend fewer joules per bit, in amplifiers that are efficient at the peaks of a signal, in algorithms
that let enormous arrays be processed with modest energy, in the coordination of many cheap devices
instead of a few expensive ones, in learned methods for the parts of the problem that resist
modelling, and in the systems questions---energy, spectrum sharing, security---where the obstacles
are as much economic and regulatory as technical.

And look at the problems no one is working on because they seem unglamorous. The telephone network
was built by people who cared about the reliability of a relay contact; the modern internet rests on a
line code that two IBM engineers designed for a mainframe storage channel (Chapter~\ref{ch:baseband}). Most of what matters is
detail, done well.

\begin{keyidea}[The apprentice's inheritance]
You do not need to remember every standard in this book; standards change. You need to be able to
look at any system, past or future, and ask: what is the bandwidth, what is the SNR, where does the
noise come from, what does the channel do, how many degrees of freedom are there, what is the receiver
matched to, and what does it cost in power, silicon and money? If you can answer those, you can
understand 6G, and 7G, and whatever comes after the generations stop being counted. The rest is
practice. Go and build something.
\end{keyidea}

%======================================================================
\section{Summary}
\label{sec:ch25:summary}
%======================================================================

\begin{itemize}
\item 6G is being defined now: ITU-R Recommendation M.2160 (2023) sets the IMT-2030 framework with six
usage scenarios, including three new ones (ubiquitous connectivity, AI and communication, integrated
sensing and communication), four overarching aspects and fifteen capabilities; 3GPP studies 6G in
Release~20 and plans the first specifications in Release~21, for deployments around 2030. The example
targets ask for modest improvement factors over IMT-2020; the new capabilities are the real change.
\item The upper mid-band (7--24\,GHz) is the most likely new capacity layer: with the same panel size,
array gain offsets the extra path loss. Sub-THz bands offer tens of gigahertz but are limited by
directivity, blockage and hardware: PA power, phase noise (which scales as $20\log_{10}K$ with frequency
multiplication) and converter power (the Walden FOM).
\item OTFS places data in the delay--Doppler domain via the ISFFT and an OFDM modulator, turning a
doubly dispersive channel into a sparse 2-D convolution with full-frame diversity and compact channel
estimation; AFDM achieves similar separability with chirps and a single DAFT. OFDM remains the
baseline; delay--Doppler ideas are likely to enter as options, estimators and sensing tools.
\item Large arrays put users in the near field ($r<2D^2/\lambda$), where beams become focused spots with
finite depth and users can be separated by range. Cell-free massive MIMO removes the cell edge but
needs fronthaul, synchronisation and user-centric clustering to scale. RIS obey the product-distance
law: they must be large, and they belong near one end of the link.
\item ISAC turns the network into a radar: range resolution $c/2B$, velocity resolution $\lambda/2T_{\text{obs}}$,
with a fundamental trade-off between deterministic signals that sense well and random signals that
carry information.
\item AI enters the air interface at three levels: inside one vendor's box, across the interface (CSI
compression, beam prediction, positioning; 3GPP Releases~18--19) and end to end. Learned components
should beat the best model-based baseline and come with a fallback.
\item Satellites and HAPS join a three-dimensional network; semantic communication should be judged
against the separation theorem; energy efficiency is mostly about idle power; zero-energy devices
use backscatter; security must migrate to post-quantum algorithms; QKD is limited by loss; optical
wireless offers huge unlicensed bandwidth indoors.
\item The tools of this book---Shannon, the matched filter, the link budget, Fourier and counting
degrees of freedom---are how to judge every claim about 6G and its successors.
\end{itemize}

\begin{labbox}[Frontier experiments]
Two notebooks accompany this chapter:
\begin{itemize}
\item \lab{moderncomms-labs/labs/lab12\_frontiers.py}
\item \lab{moderncomms-labs/labs/lab32\_6g\_waveforms.py} (planned)
\end{itemize}
The first contains four self-contained experiments: OTFS versus OFDM on a four-path doubly dispersive
channel with a $16\times16$ grid and LMMSE detection, including a plot of the sparse delay--Doppler
effective channel (Section~\ref{sec:ch25:waveforms}); an interactive OFDM radar at 28\,GHz with NR
numerology~3 that draws range--Doppler maps for two adjustable targets
(Section~\ref{sec:ch25:isac}); an interactive RIS experiment showing the $N^2$ law and the cost of
1- to 4-bit phase quantisation (Section~\ref{sec:ch25:mimo}); and a small neural-network demapper,
written in NumPy, that learns 16-QAM decision regions through a saturating amplifier with phase
noise (Section~\ref{sec:ch25:ai}). Its exercises include adding an ICI-aware OFDM equaliser and
implementing AFDM. The planned second lab extends these with the tools of this chapter's figure
script, using the new \lab{commlib.sixg} module: OTFS, AFDM and OFDM side by side with fractional
Doppler and embedded-pilot channel estimation; near-field beamfocusing maps and depth-of-focus
measurements; an ISAC range--Doppler processor with constellation-dependent sidelobe floors; and the
RIS product-distance law with geometry sliders. The figure script
\lab{book/figscripts/ch25\_figs.py} reproduces every data figure in this chapter, including the
autoencoder of Figure~\ref{fig:ch25_autoencoder}.
\end{labbox}

\problems
\begin{enumerate}[label=\thechapter.\arabic*]
\item Using Table~\ref{tab:ch25:kpi}, compute the improvement factor of IMT-2030 over IMT-2020 for each
capability. Which require new spectrum, which new antennas, and which new protocols? Argue your answers
from $C=B\log_2(1+\mathrm{SNR})$ and from the latency budget of an air interface.
\item A 5G panel at 3.5\,GHz has 192 elements at half-wavelength spacing. How many elements fit in the
same area at 7, 10 and 15\,GHz? For a fixed-gain handset antenna, compare the downlink received power
at the three frequencies. What changes on the uplink?
\item Repeat the 300\,GHz link budget for a D-band link at 140\,GHz over 200\,m with 10\,dBm transmit power,
35\,dBi antennas at each end, a 10\,GHz channel and an 8\,dB noise figure, including gaseous absorption
from Figure~\ref{fig:ch25_absorption} and 15\,dB/km of rain. What spectral efficiency is supported?
\item An oscillator at 10\,GHz has phase noise $-110$\,dBc/Hz at 1\,MHz offset. Estimate the phase noise
at the same offset if it is multiplied to 140\,GHz. If the phase-noise spectrum is flat to 1\,MHz and
falls at 20\,dB/decade beyond, estimate the RMS phase error and the ICI-limited SNR for subcarrier
spacings of 120\,kHz and 3.84\,MHz.
\item Using the Walden FOM~\eqref{eq:ch25:walden} with 30\,fJ/step, compare the ADC power of (a) 64
digital chains at 500\,MS/s with 10 ENOB and (b) 256 chains at 4\,GS/s with 4 ENOB. Which has more
total ``bits per second'' of conversion?
\item Show that with rectangular pulses the ISFFT followed by OFDM modulation of each column gives
$\vect{s}=\operatorname{vec}(\vect{x}\vect{F}_N^{\mathsf H})$, equation~\eqref{eq:ch25:otfsvec}. Then
show that a path $\bm{\Delta}^k\bm{\Pi}^\ell$ becomes a 2-D cyclic shift in the delay--Doppler
domain, and find the phase factor.
\item Design an OTFS grid for a LEO satellite link at 2\,GHz with a residual Doppler of $\pm2$\,kHz after
pre-compensation, a delay spread of 1\,$\mu$s and a latency budget of 1\,ms per frame. Choose
$\Delta f$, $M$ and $N$, and estimate the embedded-pilot overhead.
\item Show that AFDM with $c_1=c_2=0$ is OFDM. For $N=256$, a maximum Doppler of 3 bins and a maximum delay
of 10 samples, find $c_1$ from~\eqref{eq:ch25:c1} with $\xi=1$ and check the full-diversity condition.
\item Derive the Rayleigh distance~\eqref{eq:ch25:rayleigh} from the Fresnel
approximation~\eqref{eq:ch25:fresnel}. Then compute $d_R$ for a 0.5\,m array at 7, 15 and 140\,GHz and
for a 4\,m XL-MIMO facade at 3.5\,GHz.
\item An RIS of $0.5\,\text{m}\times0.5\,\text{m}$ with half-wavelength elements at 28\,GHz is placed 5\,m
from a user and 50\,m from a base station; the direct path is 52\,m long. How many elements does it
have? Compare the RIS path gain with the direct path using~\eqref{eq:ch25:productlaw}. How much blockage
of the direct path makes the RIS worthwhile? Repeat with 1-bit phase control.
\item Using~\eqref{eq:ch25:radarres}, find the range and velocity resolution and the unambiguous
range and velocity of an ISAC system at 7.5\,GHz with 30\,kHz subcarriers, 3264 subcarriers and 280
symbols. What limits the useful range if the normal cyclic prefix is used?
\item Show that the data-induced sidelobe floor of a correlation OFDM radar is $(\kappa-1)/M$,
equation~\eqref{eq:ch25:floor}, and compute $\kappa$ for 16-QAM, 64-QAM and 16-APSK with ring ratio 2.85.
\item Using the EARTH model~\eqref{eq:ch25:earth}, find the daily energy of a macro site whose load is
70\% for 14 hours and 10\% for 10 hours, with and without micro-sleep. What would deep sleep at 20\,W per
chain for the night hours save?
\item In BB84, an eavesdropper intercepts a fraction $f$ of the photons and resends them in a random
basis. Show that the QBER she introduces is $f/4$. What fraction can she intercept before the 11\%
threshold is reached? Compute the PLOB bound on the key rate at 50, 150 and 300\,km of fibre.
\item \emph{(Simulation)} Extend \texttt{fig\_otfs\_ber} in \texttt{ch25\_figs.py}: add a rate-1/2
convolutional code with interleaving across the frame to the OFDM system (use the Viterbi decoder from
\texttt{commlib}) and compare coded OFDM with uncoded and coded OTFS. How much of the gap remains?
\item \emph{(Simulation)} Reproduce Figure~\ref{fig:ch25_nearfield}. Then place two users at the same
angle, 4\,m and 12\,m from the array, and compute the SINR of each with far-field ZF beamforming and
with near-field ZF beamforming. Explain the difference.
\item \emph{(Simulation)} Train the autoencoder of Figure~\ref{fig:ch25_autoencoder} for $M=16$ through
a Rapp power amplifier with smoothness $p=2$ and an output backoff of 1\,dB. Compare the learned
constellation with 16-QAM and 16-APSK.
\end{enumerate}

\furtherreading
\begin{itemize}
\item Recommendation ITU-R M.2160-0, ``Framework and overall objectives of the future development of
IMT for 2030 and beyond,'' ITU, 2023. Short and readable; the source of the usage scenarios and
capabilities. Compare with Report ITU-R M.2410 for the IMT-2020 requirements.
\item T.~S.~Rappaport \emph{et al.}, ``Wireless communications and applications above 100\,GHz:
opportunities and challenges for 6G and beyond,'' \emph{IEEE Access}, vol.~7, 2019. A broad survey of
sub-THz propagation, hardware and applications from the group that made the case for mmWave 5G.
\item R.~Hadani \emph{et al.}, ``Orthogonal time frequency space modulation,'' \emph{Proc. IEEE WCNC},
2017; and P.~Raviteja, K.~T.~Phan and Y.~Hong, ``Embedded pilot-aided channel estimation for OTFS in
delay--Doppler channels,'' \emph{IEEE Trans. Vehicular Technology}, vol.~68, no.~5, 2019. The original
OTFS paper and the most practical channel-estimation scheme.
\item A.~Bemani, N.~Ksairi and M.~Kountouris, ``Affine frequency division multiplexing for next
generation wireless communications,'' \emph{IEEE Trans. Wireless Communications}, vol.~22, no.~11,
2023. AFDM, its diversity analysis and the choice of chirp parameters.
\item E.~Bj\"ornson, L.~Sanguinetti, H.~Wymeersch, J.~Hoydis and T.~L.~Marzetta, ``Massive MIMO is a
reality---What is next? Five promising research directions for antenna arrays,'' \emph{Digital Signal
Processing}, vol.~94, 2019. A clear, sober overview of cell-free MIMO, XL-MIMO, RIS and more.
\item \"O.~T.~Demir, E.~Bj\"ornson and L.~Sanguinetti, \emph{Foundations of User-Centric Cell-Free
Massive MIMO}, Foundations and Trends in Signal Processing, Now Publishers, 2021. The rigorous treatment
of cell-free processing levels, scalability and fronthaul, with code.
\item E.~Bj\"ornson, \"O.~\"Ozdogan and E.~G.~Larsson, ``Intelligent reflecting surface versus
decode-and-forward: how large surfaces are needed to beat relaying?'' \emph{IEEE Wireless
Communications Letters}, vol.~9, no.~2, 2020. A short paper every RIS enthusiast should read.
\item F.~Liu \emph{et al.}, ``Integrated sensing and communications: toward dual-functional wireless
networks for 6G and beyond,'' \emph{IEEE J. Selected Areas in Communications}, vol.~40, no.~6, 2022;
and Y.~Xiong \emph{et al.}, ``On the fundamental tradeoff of integrated sensing and communications
under Gaussian channels,'' \emph{IEEE Trans. Information Theory}, vol.~69, no.~9, 2023. The ISAC survey
and the deterministic--random trade-off.
\item T.~O'Shea and J.~Hoydis, ``An introduction to deep learning for the physical layer,''
\emph{IEEE Trans. Cognitive Communications and Networking}, vol.~3, no.~4, 2017. The paper that launched
end-to-end learning; read it together with 3GPP TR~38.843 for the standards view of AI/ML in the air
interface.
\item D.~G\"und\"uz \emph{et al.}, ``Beyond transmitting bits: context, semantics, and task-oriented
communications,'' \emph{IEEE J. Selected Areas in Communications}, vol.~41, no.~1, 2023. A balanced
survey of semantic and goal-oriented communication by researchers who take the theory seriously.
\item R.~H.~Walden, ``Analog-to-digital converter survey and analysis,'' \emph{IEEE J. Selected Areas in
Communications}, vol.~17, no.~4, 1999; and B.~Murmann, ``ADC Performance Survey,'' online, updated
annually. Where the figure of merit comes from, and where converters actually are.
\item C.~H.~Bennett and G.~Brassard, ``Quantum cryptography: public key distribution and coin tossing,''
\emph{Proc. IEEE Int. Conf. Computers, Systems and Signal Processing}, Bangalore, 1984; and NIST FIPS~203,
204 and 205 (2024) for the post-quantum standards. Two routes to keys that survive a quantum computer.
\end{itemize}
